\documentclass[../main.tex]{subfiles}
\begin{document}

\chapter{Virtualization}
\lessoninfo{
  video={P3L6},
  textbook={OSTEP \cite{ostep} app.~B; Silberschatz et al.\ \cite{silberschatz2018} ch.~18},
  keywords={virtualization, virtual machine, hypervisor, bare-metal hypervisor, hosted hypervisor, protection ring, trap-and-emulate, binary translation, paravirtualization, hypercall, shadow page table, split-device driver model, hardware-assisted virtualization, extended page table},
}

The preceding lessons showed how an operating system manages the CPUs, the
memory and the devices of a machine on behalf of the processes that run on
it.  Virtualization moves this role one level down: a thin layer of
software manages the hardware on behalf of several complete operating
systems, each of which behaves as if it owned the machine.  This lesson
defines virtualization, explains why it has become widespread, and compares
the two ways in which virtualization software is organized.  It then
examines how the CPU, memory and devices are virtualized, first with
software techniques alone and then with the support that current processors
provide.  The protection boundary of Lesson~1, the page tables of Lesson~8
and the device drivers of Lesson~11 reappear throughout, now with one more
layer of software beneath the operating system.

\section{What is virtualization?}
\label{sec:l12-what}

Virtualization is an old idea.  It originated at IBM in the 1960s, when
computing was done on a small number of large and expensive mainframes, and
the owner of such a machine wanted to use it for workloads that had been
developed for different operating systems.\source{P3L6, what is
virtualization, defining virtualization; Popek and Goldberg
\cite{popek1974}; OSTEP app.~B; Silberschatz et al.\ ch.~18.}

\begin{definition}[Virtualization]
  Allowing several operating systems, each together with its applications,
  to execute concurrently on the same physical machine is called
  \term{virtualization}[仮想化].
\end{definition}

None of these operating systems is in control of the hardware any more.
Each runs in an environment of its own and is given the illusion that it
owns the hardware resources---CPUs, memory and devices---although the
resources it sees are \emph{virtual resources} that stand in for a share of
the physical ones (\cref{fig:l12-vmm}).

\begin{definition}[Virtual machine]
  A \term{virtual machine}[仮想マシン] (VM) is an operating system together
  with its applications and the virtual resources on which they run.
\end{definition}
\index{VM|see{virtual machine}}

\begin{figure}[htbp]
  \centering
  % Virtualization: several virtual machines, each with its own guest OS,
% applications and virtual resources, on one virtual machine monitor.
% Usage: % Virtualization: several virtual machines, each with its own guest OS,
% applications and virtual resources, on one virtual machine monitor.
% Usage: % Virtualization: several virtual machines, each with its own guest OS,
% applications and virtual resources, on one virtual machine monitor.
% Usage: \input{figures/l12-vmm}   (width ~96 mm)
\begin{tikzpicture}[x=1mm, y=1mm, font=\footnotesize\sffamily,
  app/.style={draw, rounded corners=1pt, fill=white, minimum height=5mm,
              minimum width=11.5mm, inner sep=0.5pt, font=\scriptsize\sffamily},
  gos/.style={draw, fill=ln-accent!12, minimum height=6mm,
              minimum width=26mm, inner sep=0.5pt},
  vres/.style={draw, densely dashed, fill=white, minimum height=5mm,
               minimum width=26mm, inner sep=0.5pt, font=\scriptsize\sffamily},
  vm/.style={draw=ln-rule, rounded corners=2pt, fill=ln-box},
  layer/.style={draw, minimum height=7mm, minimum width=96mm, inner sep=1pt},
  ttl/.style={font=\footnotesize\sffamily\bfseries}]
  \node[layer, fill=white] at (48,3.5)
    {\iflangja{物理ハードウェア（CPU，メモリ，デバイス）}{physical hardware (CPUs, memory, devices)}};
  \node[layer, fill=ln-accent!22] at (48,12.5)
    {\iflangja{仮想マシンモニタ（ハイパーバイザ）}{virtual machine monitor (hypervisor)}};
  \foreach \x/\n in {16/1, 48/2, 80/3} {
    \draw[vm] (\x-15,18) rectangle (\x+15,48);
    \node[ttl] at (\x,44.5) {VM \n};
    \node[app] at (\x-7,37.5) {\iflangja{アプリ}{app}};
    \node[app] at (\x+7,37.5) {\iflangja{アプリ}{app}};
    \node[gos] at (\x,29.5) {\iflangja{ゲスト OS}{guest OS}};
    \node[vres] at (\x,22.5) {\iflangja{仮想資源}{virtual resources}};
  }
\end{tikzpicture}
   (width ~96 mm)
\begin{tikzpicture}[x=1mm, y=1mm, font=\footnotesize\sffamily,
  app/.style={draw, rounded corners=1pt, fill=white, minimum height=5mm,
              minimum width=11.5mm, inner sep=0.5pt, font=\scriptsize\sffamily},
  gos/.style={draw, fill=ln-accent!12, minimum height=6mm,
              minimum width=26mm, inner sep=0.5pt},
  vres/.style={draw, densely dashed, fill=white, minimum height=5mm,
               minimum width=26mm, inner sep=0.5pt, font=\scriptsize\sffamily},
  vm/.style={draw=ln-rule, rounded corners=2pt, fill=ln-box},
  layer/.style={draw, minimum height=7mm, minimum width=96mm, inner sep=1pt},
  ttl/.style={font=\footnotesize\sffamily\bfseries}]
  \node[layer, fill=white] at (48,3.5)
    {\iflangja{物理ハードウェア（CPU，メモリ，デバイス）}{physical hardware (CPUs, memory, devices)}};
  \node[layer, fill=ln-accent!22] at (48,12.5)
    {\iflangja{仮想マシンモニタ（ハイパーバイザ）}{virtual machine monitor (hypervisor)}};
  \foreach \x/\n in {16/1, 48/2, 80/3} {
    \draw[vm] (\x-15,18) rectangle (\x+15,48);
    \node[ttl] at (\x,44.5) {VM \n};
    \node[app] at (\x-7,37.5) {\iflangja{アプリ}{app}};
    \node[app] at (\x+7,37.5) {\iflangja{アプリ}{app}};
    \node[gos] at (\x,29.5) {\iflangja{ゲスト OS}{guest OS}};
    \node[vres] at (\x,22.5) {\iflangja{仮想資源}{virtual resources}};
  }
\end{tikzpicture}
   (width ~96 mm)
\begin{tikzpicture}[x=1mm, y=1mm, font=\footnotesize\sffamily,
  app/.style={draw, rounded corners=1pt, fill=white, minimum height=5mm,
              minimum width=11.5mm, inner sep=0.5pt, font=\scriptsize\sffamily},
  gos/.style={draw, fill=ln-accent!12, minimum height=6mm,
              minimum width=26mm, inner sep=0.5pt},
  vres/.style={draw, densely dashed, fill=white, minimum height=5mm,
               minimum width=26mm, inner sep=0.5pt, font=\scriptsize\sffamily},
  vm/.style={draw=ln-rule, rounded corners=2pt, fill=ln-box},
  layer/.style={draw, minimum height=7mm, minimum width=96mm, inner sep=1pt},
  ttl/.style={font=\footnotesize\sffamily\bfseries}]
  \node[layer, fill=white] at (48,3.5)
    {\iflangja{物理ハードウェア（CPU，メモリ，デバイス）}{physical hardware (CPUs, memory, devices)}};
  \node[layer, fill=ln-accent!22] at (48,12.5)
    {\iflangja{仮想マシンモニタ（ハイパーバイザ）}{virtual machine monitor (hypervisor)}};
  \foreach \x/\n in {16/1, 48/2, 80/3} {
    \draw[vm] (\x-15,18) rectangle (\x+15,48);
    \node[ttl] at (\x,44.5) {VM \n};
    \node[app] at (\x-7,37.5) {\iflangja{アプリ}{app}};
    \node[app] at (\x+7,37.5) {\iflangja{アプリ}{app}};
    \node[gos] at (\x,29.5) {\iflangja{ゲスト OS}{guest OS}};
    \node[vres] at (\x,22.5) {\iflangja{仮想資源}{virtual resources}};
  }
\end{tikzpicture}

  \caption{Virtualization.  Each VM consists of a guest OS, its
    applications and the virtual resources presented to it; the virtual
    machine monitor manages the physical hardware beneath all VMs.}
  \label{fig:l12-vmm}
\end{figure}

\begin{definition}[Virtual machine monitor]
  The \term{virtual machine monitor}[仮想マシンモニタ] (VMM), also called
  the \emph{hypervisor}, is the virtualization layer beneath the VMs.  It
  manages the physical hardware, provides the virtual resources to the VMs,
  coordinates their access to the physical resources, and isolates them
  from one another.
\end{definition}
\index{VMM|see{virtual machine monitor}}%
\index{hypervisor|see{virtual machine monitor}}

The operating system that runs inside a VM is called a
\term{guest operating system}[ゲストOS] (guest OS), and a VM as a whole is
also called a \emph{guest domain}.

Popek and Goldberg \cite{popek1974} gave the classic definition: a virtual
machine is \enquote{an efficient, isolated duplicate of the real machine}.
The VMM that supports it must provide three properties.%
\margin{Popek and Goldberg name the three properties equivalence,
efficiency and resource control.}
\begin{description}
  \item[Fidelity] Programs, operating systems included, run in an
    environment essentially identical to the original machine and behave
    as they would on it, apart from timing and the amount of available
    resources.  This \term{fidelity}[忠実性] is what lets an unmodified
    operating system run in a VM; paravirtualization
    (\cref{sec:l12-pv}) deliberately relaxes it.
  \item[Performance] Programs in a VM show at worst a minor decrease in
    speed.  Most instructions must therefore execute directly on the
    hardware rather than be interpreted by the VMM.
  \item[Safety and isolation] The VMM is in complete control of the system
    resources.  A VM cannot access resources that the VMM has not allocated
    to it, and it cannot interfere with other VMs or with the VMM itself.
\end{description}

\begin{note}
  The definition excludes technologies that are also called virtual
  machines.  A high-level language virtual machine executes programs
  compiled for an abstract instruction set, not for the real machine, so
  it is not a duplicate of the hardware beneath it.  An emulator that
  implements a different hardware platform in software duplicates that
  platform, not the real machine, and it interprets every instruction, so
  it does not meet the performance requirement.
\end{note}

\section{Benefits of virtualization}
\label{sec:l12-benefits}

Running several operating systems on one physical machine has a number of
benefits.\source{P3L6, benefits of virtualization; Rosenblum and Garfinkel
\cite{rosenblum2005}; Silberschatz et al.\ ch.~18.}
\begin{description}
  \item[Consolidation] The workloads of several physical machines, each
    with its own operating system and applications, can run as VMs on fewer
    machines.  This \term{consolidation}[サーバ統合] decreases cost---fewer
    machines to buy, house, power and cool---and improves manageability,
    because fewer machines need administration.
  \item[Migration] A VM is decoupled from the physical machine, so its
    complete state, operating system and applications included, can be
    moved from one machine to another, for example before a machine is
    taken down for maintenance.  Migration improves availability and
    reliability.
  \item[Security] Bugs and malicious behaviour are contained within the VM
    in which they occur and do not affect the other VMs.
  \item[Debugging] New operating systems and new operating system features
    can be brought up and tested quickly in a VM, without a dedicated
    machine; a crash takes down only that VM.
  \item[Support for legacy operating systems] An operating system that no
    longer runs on current hardware, together with the applications that
    depend on it, can continue to run in a VM next to current systems.
\end{description}

Although IBM mainframes supported virtualization from the 1960s on, it
remained confined to them until the late 1990s \cite{rosenblum2005}.
Mainframes were expensive and few.  Once commodity hardware had become cheap and its
operating systems supported multitasking, it was simpler to buy another
machine for another application than to share one machine among several
operating systems, and interest in VMMs faded.  This practice eventually
created the problems that revived them.  Data centers filled with machines
that each ran a single application and were underutilized; their number
grew, and with it the number of administrators needed and the cost of
power and cooling.  Consolidating the applications onto fewer machines
addresses all these costs at once, but it requires virtualizing commodity
x86 hardware, which, as \cref{sec:l12-x86} shows, had not been designed for
it.

\begin{example}
  A company runs 24 servers, each dedicated to one application and
  averaging \SI{8}{\percent} CPU utilization.  Their combined load
  corresponds to $24 \times 0.08 = 1.92$ fully used servers.  If each
  physical machine is loaded to at most \SI{70}{\percent}, leaving headroom
  for load peaks and for the overhead of virtualization, and memory and
  I/O capacity suffice, the 24 workloads fit as VMs on
  $\lceil 1.92/0.7 \rceil = 3$ machines, each at about \SI{64}{\percent}
  utilization.  The number of machines to power, cool and administer drops
  by a factor of eight.
\end{example}

\section{Virtualization models}
\label{sec:l12-models}

Virtualization software is organized in one of two ways, which differ in
what runs directly on the hardware and, in particular, in who owns the
device drivers.\source{P3L6, virtualization models: bare-metal, hosted;
Barham et al.\ \cite{barham2003}; Kivity et al.\ \cite{kivity2007};
Silberschatz et al.\ ch.~18.}

\subsection{Bare-metal hypervisors}

\begin{definition}[Bare-metal hypervisor]
  A \term{bare-metal hypervisor}[ベアメタル型ハイパーバイザ], also called a
  hypervisor-based or \emph{type~1} VMM, runs directly on the hardware.  It
  manages all hardware resources and supports the execution of the VMs.
\end{definition}

Such a hypervisor must also control the devices.  Devices are diverse, and
supporting all of them would put a large body of device drivers into the
hypervisor.  Bare-metal designs therefore often run a privileged
\emph{service VM} next to the guest VMs (\cref{fig:l12-models}a).  The
service VM runs a standard operating system with the device drivers, has
direct access to the devices, and performs configuration and management
tasks, so that the hypervisor itself stays small.

\begin{figure}[htbp]
  \centering
  % Virtualization models: (a) bare-metal hypervisor with a privileged
% service VM that owns the device drivers (Xen: dom0 and domU guests);
% (b) hosted hypervisor: a module in the host OS (KVM in Linux), each VM a
% host process with QEMU emulating its devices.
% Usage: % Virtualization models: (a) bare-metal hypervisor with a privileged
% service VM that owns the device drivers (Xen: dom0 and domU guests);
% (b) hosted hypervisor: a module in the host OS (KVM in Linux), each VM a
% host process with QEMU emulating its devices.
% Usage: % Virtualization models: (a) bare-metal hypervisor with a privileged
% service VM that owns the device drivers (Xen: dom0 and domU guests);
% (b) hosted hypervisor: a module in the host OS (KVM in Linux), each VM a
% host process with QEMU emulating its devices.
% Usage: \input{figures/l12-models}   (width ~116 mm)
\begin{tikzpicture}[x=1mm, y=1mm, font=\footnotesize\sffamily,
  >={Stealth[length=1.6mm]},
  box/.style={draw, fill=white, inner sep=0.5pt, align=center,
              font=\scriptsize\sffamily},
  gos/.style={box, fill=ln-accent!12},
  vm/.style={draw=ln-rule, rounded corners=2pt, fill=ln-box},
  svm/.style={draw=ln-accent, rounded corners=2pt, fill=ln-accent!8},
  lab/.style={font=\scriptsize\sffamily, align=center},
  hdr/.style={font=\scriptsize\sffamily\bfseries, align=center},
  pcap/.style={font=\footnotesize\sffamily}]
  % ================= (a) bare-metal
  \draw[fill=white] (0,0) rectangle (55,7)
    node[midway] {\iflangja{ハードウェア}{hardware}};
  \draw[fill=ln-accent!22] (0,9) rectangle (55,16)
    node[midway] {\iflangja{ハイパーバイザ（VMM）}{hypervisor (VMM)}};
  % service VM
  \draw[svm] (0,18) rectangle (19,52);
  \node[hdr] at (9.5,49) {\iflangja{サービス VM}{service VM}};
  \node[lab] at (9.5,45.5) {(dom0)};
  \node[box, minimum width=16mm, minimum height=7mm] at (9.5,36.5)
    {\iflangja{管理ツール}{management}};
  \node[gos, minimum width=16mm, minimum height=9mm] (drv) at (9.5,25.5)
    {OS +\\\iflangja{ドライバ}{drivers}};
  \draw[<->, thick, densely dashed, ln-accent] (drv.south) -- (9.5,7);
  % guest VMs
  \foreach \x in {29, 47} {
    \draw[vm] (\x-8,18) rectangle (\x+8,52);
    \node[hdr] at (\x,49) {\iflangja{ゲスト VM}{guest VM}};
    \node[lab] at (\x,45.5) {(domU)};
    \node[box, rounded corners=1pt, minimum width=13mm, minimum height=6mm]
      at (\x,36.5) {\iflangja{アプリ}{apps}};
    \node[gos, minimum width=13mm, minimum height=7mm] at (\x,26)
      {\iflangja{ゲスト OS}{guest OS}};
  }
  \node[pcap] at (27.5,-4) {\iflangja{(a) ベアメタル型}{(a) bare-metal}};
  % ================= (b) hosted
  \begin{scope}[xshift=61mm]
  \draw[fill=white] (0,0) rectangle (55,7)
    node[midway] {\iflangja{ハードウェア}{hardware}};
  \draw[fill=ln-box] (0,9) rectangle (55,26);
  \node[hdr] at (27.5,22.5) {\iflangja{ホスト OS（Linux）}{host OS (Linux)}};
  \node[box, minimum width=24mm, minimum height=7mm] at (14,14.5)
    {\iflangja{デバイスドライバ}{device drivers}};
  \node[box, fill=ln-accent!22, minimum width=24mm, minimum height=7mm]
    at (41,14.5) {\iflangja{KVM モジュール}{KVM module}};
  % native process
  \draw[fill=white, rounded corners=2pt] (0,28) rectangle (13,52)
    node[midway, lab] {\iflangja{通常の\\プロセス}{native\\process}};
  % VMs as host processes
  \foreach \x in {24.5, 45.5} {
    \draw[vm] (\x-9.5,28) rectangle (\x+9.5,52);
    \node[hdr] at (\x,49.3) {VM};
    \node[box, rounded corners=1pt, minimum width=15mm, minimum height=5mm]
      at (\x,43.5) {\iflangja{アプリ}{apps}};
    \node[gos, minimum width=16mm, minimum height=5.5mm] at (\x,37)
      {\iflangja{ゲスト OS}{guest OS}};
    \node[box, densely dashed, minimum width=16mm, minimum height=5mm]
      at (\x,31) {QEMU};
  }
  \node[pcap] at (27.5,-4) {\iflangja{(b) ホスト型}{(b) hosted}};
  \end{scope}
\end{tikzpicture}
   (width ~116 mm)
\begin{tikzpicture}[x=1mm, y=1mm, font=\footnotesize\sffamily,
  >={Stealth[length=1.6mm]},
  box/.style={draw, fill=white, inner sep=0.5pt, align=center,
              font=\scriptsize\sffamily},
  gos/.style={box, fill=ln-accent!12},
  vm/.style={draw=ln-rule, rounded corners=2pt, fill=ln-box},
  svm/.style={draw=ln-accent, rounded corners=2pt, fill=ln-accent!8},
  lab/.style={font=\scriptsize\sffamily, align=center},
  hdr/.style={font=\scriptsize\sffamily\bfseries, align=center},
  pcap/.style={font=\footnotesize\sffamily}]
  % ================= (a) bare-metal
  \draw[fill=white] (0,0) rectangle (55,7)
    node[midway] {\iflangja{ハードウェア}{hardware}};
  \draw[fill=ln-accent!22] (0,9) rectangle (55,16)
    node[midway] {\iflangja{ハイパーバイザ（VMM）}{hypervisor (VMM)}};
  % service VM
  \draw[svm] (0,18) rectangle (19,52);
  \node[hdr] at (9.5,49) {\iflangja{サービス VM}{service VM}};
  \node[lab] at (9.5,45.5) {(dom0)};
  \node[box, minimum width=16mm, minimum height=7mm] at (9.5,36.5)
    {\iflangja{管理ツール}{management}};
  \node[gos, minimum width=16mm, minimum height=9mm] (drv) at (9.5,25.5)
    {OS +\\\iflangja{ドライバ}{drivers}};
  \draw[<->, thick, densely dashed, ln-accent] (drv.south) -- (9.5,7);
  % guest VMs
  \foreach \x in {29, 47} {
    \draw[vm] (\x-8,18) rectangle (\x+8,52);
    \node[hdr] at (\x,49) {\iflangja{ゲスト VM}{guest VM}};
    \node[lab] at (\x,45.5) {(domU)};
    \node[box, rounded corners=1pt, minimum width=13mm, minimum height=6mm]
      at (\x,36.5) {\iflangja{アプリ}{apps}};
    \node[gos, minimum width=13mm, minimum height=7mm] at (\x,26)
      {\iflangja{ゲスト OS}{guest OS}};
  }
  \node[pcap] at (27.5,-4) {\iflangja{(a) ベアメタル型}{(a) bare-metal}};
  % ================= (b) hosted
  \begin{scope}[xshift=61mm]
  \draw[fill=white] (0,0) rectangle (55,7)
    node[midway] {\iflangja{ハードウェア}{hardware}};
  \draw[fill=ln-box] (0,9) rectangle (55,26);
  \node[hdr] at (27.5,22.5) {\iflangja{ホスト OS（Linux）}{host OS (Linux)}};
  \node[box, minimum width=24mm, minimum height=7mm] at (14,14.5)
    {\iflangja{デバイスドライバ}{device drivers}};
  \node[box, fill=ln-accent!22, minimum width=24mm, minimum height=7mm]
    at (41,14.5) {\iflangja{KVM モジュール}{KVM module}};
  % native process
  \draw[fill=white, rounded corners=2pt] (0,28) rectangle (13,52)
    node[midway, lab] {\iflangja{通常の\\プロセス}{native\\process}};
  % VMs as host processes
  \foreach \x in {24.5, 45.5} {
    \draw[vm] (\x-9.5,28) rectangle (\x+9.5,52);
    \node[hdr] at (\x,49.3) {VM};
    \node[box, rounded corners=1pt, minimum width=15mm, minimum height=5mm]
      at (\x,43.5) {\iflangja{アプリ}{apps}};
    \node[gos, minimum width=16mm, minimum height=5.5mm] at (\x,37)
      {\iflangja{ゲスト OS}{guest OS}};
    \node[box, densely dashed, minimum width=16mm, minimum height=5mm]
      at (\x,31) {QEMU};
  }
  \node[pcap] at (27.5,-4) {\iflangja{(b) ホスト型}{(b) hosted}};
  \end{scope}
\end{tikzpicture}
   (width ~116 mm)
\begin{tikzpicture}[x=1mm, y=1mm, font=\footnotesize\sffamily,
  >={Stealth[length=1.6mm]},
  box/.style={draw, fill=white, inner sep=0.5pt, align=center,
              font=\scriptsize\sffamily},
  gos/.style={box, fill=ln-accent!12},
  vm/.style={draw=ln-rule, rounded corners=2pt, fill=ln-box},
  svm/.style={draw=ln-accent, rounded corners=2pt, fill=ln-accent!8},
  lab/.style={font=\scriptsize\sffamily, align=center},
  hdr/.style={font=\scriptsize\sffamily\bfseries, align=center},
  pcap/.style={font=\footnotesize\sffamily}]
  % ================= (a) bare-metal
  \draw[fill=white] (0,0) rectangle (55,7)
    node[midway] {\iflangja{ハードウェア}{hardware}};
  \draw[fill=ln-accent!22] (0,9) rectangle (55,16)
    node[midway] {\iflangja{ハイパーバイザ（VMM）}{hypervisor (VMM)}};
  % service VM
  \draw[svm] (0,18) rectangle (19,52);
  \node[hdr] at (9.5,49) {\iflangja{サービス VM}{service VM}};
  \node[lab] at (9.5,45.5) {(dom0)};
  \node[box, minimum width=16mm, minimum height=7mm] at (9.5,36.5)
    {\iflangja{管理ツール}{management}};
  \node[gos, minimum width=16mm, minimum height=9mm] (drv) at (9.5,25.5)
    {OS +\\\iflangja{ドライバ}{drivers}};
  \draw[<->, thick, densely dashed, ln-accent] (drv.south) -- (9.5,7);
  % guest VMs
  \foreach \x in {29, 47} {
    \draw[vm] (\x-8,18) rectangle (\x+8,52);
    \node[hdr] at (\x,49) {\iflangja{ゲスト VM}{guest VM}};
    \node[lab] at (\x,45.5) {(domU)};
    \node[box, rounded corners=1pt, minimum width=13mm, minimum height=6mm]
      at (\x,36.5) {\iflangja{アプリ}{apps}};
    \node[gos, minimum width=13mm, minimum height=7mm] at (\x,26)
      {\iflangja{ゲスト OS}{guest OS}};
  }
  \node[pcap] at (27.5,-4) {\iflangja{(a) ベアメタル型}{(a) bare-metal}};
  % ================= (b) hosted
  \begin{scope}[xshift=61mm]
  \draw[fill=white] (0,0) rectangle (55,7)
    node[midway] {\iflangja{ハードウェア}{hardware}};
  \draw[fill=ln-box] (0,9) rectangle (55,26);
  \node[hdr] at (27.5,22.5) {\iflangja{ホスト OS（Linux）}{host OS (Linux)}};
  \node[box, minimum width=24mm, minimum height=7mm] at (14,14.5)
    {\iflangja{デバイスドライバ}{device drivers}};
  \node[box, fill=ln-accent!22, minimum width=24mm, minimum height=7mm]
    at (41,14.5) {\iflangja{KVM モジュール}{KVM module}};
  % native process
  \draw[fill=white, rounded corners=2pt] (0,28) rectangle (13,52)
    node[midway, lab] {\iflangja{通常の\\プロセス}{native\\process}};
  % VMs as host processes
  \foreach \x in {24.5, 45.5} {
    \draw[vm] (\x-9.5,28) rectangle (\x+9.5,52);
    \node[hdr] at (\x,49.3) {VM};
    \node[box, rounded corners=1pt, minimum width=15mm, minimum height=5mm]
      at (\x,43.5) {\iflangja{アプリ}{apps}};
    \node[gos, minimum width=16mm, minimum height=5.5mm] at (\x,37)
      {\iflangja{ゲスト OS}{guest OS}};
    \node[box, densely dashed, minimum width=16mm, minimum height=5mm]
      at (\x,31) {QEMU};
  }
  \node[pcap] at (27.5,-4) {\iflangja{(b) ホスト型}{(b) hosted}};
  \end{scope}
\end{tikzpicture}

  \caption{Virtualization models.  (a) A bare-metal hypervisor with a
    privileged service VM that owns the device drivers, as in Xen.  (b) A
    hosted hypervisor: a module of the host OS, as KVM in Linux; each VM is
    a host process in which QEMU emulates the devices.}
  \label{fig:l12-models}
\end{figure}

\term{Xen}[Xen] \cite{barham2003}, available as open source and as the
commercial Citrix XenServer, follows this design.  Its VMs are called
\emph{domains}.  The privileged service VM is \term{dom0}[dom0], which runs
Linux with the device drivers and the management tools; the guest VMs are
unprivileged domains, \emph{domU}, and the Xen hypervisor itself contains
no device drivers.  VMware ESX takes the other route and places the drivers
in the hypervisor.  This is practical because ESX targets servers from a
limited number of vendors, who supply drivers for it.  ESX is administered
through remote management APIs, which replaced an earlier Linux-based
control core.

\subsection{Hosted hypervisors}

\begin{definition}[Hosted hypervisor]
  A \term{hosted hypervisor}[ホスト型ハイパーバイザ], or \emph{type~2} VMM,
  runs on top of a conventional operating system, the
  \term{host operating system}[ホストOS] (host OS), which owns all the
  hardware.  A special VMM module in the host OS provides the hardware
  interfaces to the VMs and handles the switches between VMs.
\end{definition}

A hosted hypervisor reuses what the host OS already provides.  Device
drivers, scheduling, memory management and other services need not be
re-implemented, and ordinary applications continue to run natively on the
host next to the VMs (\cref{fig:l12-models}b).

\needspace{6\baselineskip}
The \term{Kernel-based Virtual Machine}[KVM] (KVM)
\cite{kivity2007}\index{KVM|see{Kernel-based Virtual Machine}} is the
hosted hypervisor of Linux.  A loadable kernel module turns Linux into a
hypervisor.  Each VM is an ordinary Linux process, whose virtual CPUs are
threads of that process, so the Linux scheduler and memory manager serve
VMs as they serve other processes.  The devices of a VM are emulated within
its process by \term{QEMU}[QEMU], a hardware emulator running at user
level, while the instructions of the guest execute directly on the CPU with
the hardware support described in \cref{sec:l12-hw}.  Because KVM is part
of Linux, it benefits from the Linux open-source community: new features,
bug fixes and drivers for new devices become available to it without
separate effort.

\begin{keypoint}
  A bare-metal hypervisor runs directly on the hardware and typically
  leaves device drivers to a privileged service VM, such as dom0 in Xen.  A
  hosted hypervisor is a module of a host OS, such as KVM in Linux, and
  reuses the drivers and services of that OS.
\end{keypoint}

\section{Hardware protection levels}
\label{sec:l12-rings}

A guest OS is written to run in kernel mode and to execute privileged
instructions, yet under virtualization only the hypervisor may control the
hardware.\source{P3L6, hardware protection levels; Silberschatz et al.\
ch.~18.}  The two modes of Lesson~1 cannot separate three kinds of
software---applications, guest OS and hypervisor---but commodity processors
offer more than two privilege levels.

\begin{definition}[Protection ring]
  A \term{protection ring}[保護リング] is one of a hierarchy of privilege
  levels enforced by the processor; software in a ring may execute only
  the instructions and access only the resources permitted at its level.
  x86 processors provide four rings, from ring~0, the most privileged, to
  ring~3, the least privileged.
\end{definition}

Without virtualization, the OS kernel runs in ring~0 and applications in
ring~3, and rings~1 and~2 are unused (\cref{fig:l12-rings}a).  With
virtualization, the hypervisor takes ring~0, the guest OS moves to ring~1,
and applications remain in ring~3 (\cref{fig:l12-rings}b).  When the guest
OS attempts a privileged operation from ring~1, the processor traps to the
hypervisor in ring~0.

\begin{figure}[htbp]
  \centering
  % x86 protection rings: (a) native OS; (b) virtualization without hardware
% support (hypervisor in ring 0, guest OS deprivileged to ring 1);
% (c) root and non-root modes, each with four rings, VM exit / VM entry.
% Usage: % x86 protection rings: (a) native OS; (b) virtualization without hardware
% support (hypervisor in ring 0, guest OS deprivileged to ring 1);
% (c) root and non-root modes, each with four rings, VM exit / VM entry.
% Usage: % x86 protection rings: (a) native OS; (b) virtualization without hardware
% support (hypervisor in ring 0, guest OS deprivileged to ring 1);
% (c) root and non-root modes, each with four rings, VM exit / VM entry.
% Usage: \input{figures/l12-rings}   (width ~118 mm)
\begin{tikzpicture}[x=1mm, y=1mm, font=\scriptsize\sffamily,
  >={Stealth[length=1.5mm]},
  used/.style={draw, fill=ln-accent!12},
  hv/.style={draw, fill=ln-accent!30},
  free/.style={draw=ln-rule, fill=white},
  rl/.style={font=\scriptsize\sffamily, anchor=east, text=ln-muted},
  hdr/.style={font=\scriptsize\sffamily\bfseries},
  pcap/.style={font=\footnotesize\sffamily, align=center, anchor=north}]
  % row labels
  \foreach \r/\y in {3/27, 2/21, 1/15, 0/9}
    \node[rl] at (9.5,\y) {\iflangja{リング \r}{ring \r}};
  % (a) native
  \draw[used] (11,24) rectangle (31,30) node[midway] {\iflangja{アプリケーション}{applications}};
  \draw[free] (11,18) rectangle (31,24);
  \draw[free] (11,12) rectangle (31,18);
  \draw[used] (11,6)  rectangle (31,12) node[midway] {\iflangja{OS カーネル}{OS kernel}};
  \node[pcap] at (21,-1) {\iflangja{(a) 仮想化なし}{(a) native}};
  % (b) hypervisor in ring 0
  \draw[used] (34,24) rectangle (54,30) node[midway] {\iflangja{アプリケーション}{applications}};
  \draw[free] (34,18) rectangle (54,24);
  \draw[used] (34,12) rectangle (54,18) node[midway] {\iflangja{ゲスト OS}{guest OS}};
  \draw[hv]   (34,6)  rectangle (54,12) node[midway] {\iflangja{ハイパーバイザ}{hypervisor}};
  \draw[->, thick, ln-caution] (54.3,15) to[out=0, in=0, looseness=2.2]
    node[right=0.8mm, text=ln-caution] {\iflangja{}{trap}} (54.3,9);
  % The Japanese label is wider than the gap between panels (b) and (c),
  % so it is set vertically beside the arc instead of next to it.
  \iflangja{\node[text=ln-caution, rotate=90] at (60.3,12) {トラップ};}{}
  \node[pcap] at (44,-1) {\iflangja{(b) リング 0 の\\ハイパーバイザ}{(b) hypervisor\\in ring 0}};
  % (c) root / non-root
  \node[hdr] at (74,33) {\iflangja{非ルートモード}{non-root mode}};
  \node[hdr] at (108,33) {\iflangja{ルートモード}{root mode}};
  \draw[used] (64,24) rectangle (84,30) node[midway] {\iflangja{アプリケーション}{applications}};
  \draw[free] (64,18) rectangle (84,24);
  \draw[free] (64,12) rectangle (84,18);
  \draw[used] (64,6)  rectangle (84,12) node[midway] {\iflangja{ゲスト OS}{guest OS}};
  \draw[free] (98,24) rectangle (118,30);
  \draw[free] (98,18) rectangle (118,24);
  \draw[free] (98,12) rectangle (118,18);
  \draw[hv]   (98,6)  rectangle (118,12) node[midway] {\iflangja{ハイパーバイザ}{hypervisor}};
  \draw[->, thick, ln-caution] (84.3,10.2) -- (97.7,10.2);
  \draw[->, thick, ln-tip] (97.7,7.8) -- (84.3,7.8);
  \node[text=ln-caution, anchor=south] at (91,12.2) {VM exit};
  \node[text=ln-tip, anchor=north] at (91,5.8) {VM entry};
  \node[pcap] at (91,-1) {\iflangja{(c) ルートモードと非ルートモード}{(c) root and non-root modes}};
\end{tikzpicture}
   (width ~118 mm)
\begin{tikzpicture}[x=1mm, y=1mm, font=\scriptsize\sffamily,
  >={Stealth[length=1.5mm]},
  used/.style={draw, fill=ln-accent!12},
  hv/.style={draw, fill=ln-accent!30},
  free/.style={draw=ln-rule, fill=white},
  rl/.style={font=\scriptsize\sffamily, anchor=east, text=ln-muted},
  hdr/.style={font=\scriptsize\sffamily\bfseries},
  pcap/.style={font=\footnotesize\sffamily, align=center, anchor=north}]
  % row labels
  \foreach \r/\y in {3/27, 2/21, 1/15, 0/9}
    \node[rl] at (9.5,\y) {\iflangja{リング \r}{ring \r}};
  % (a) native
  \draw[used] (11,24) rectangle (31,30) node[midway] {\iflangja{アプリケーション}{applications}};
  \draw[free] (11,18) rectangle (31,24);
  \draw[free] (11,12) rectangle (31,18);
  \draw[used] (11,6)  rectangle (31,12) node[midway] {\iflangja{OS カーネル}{OS kernel}};
  \node[pcap] at (21,-1) {\iflangja{(a) 仮想化なし}{(a) native}};
  % (b) hypervisor in ring 0
  \draw[used] (34,24) rectangle (54,30) node[midway] {\iflangja{アプリケーション}{applications}};
  \draw[free] (34,18) rectangle (54,24);
  \draw[used] (34,12) rectangle (54,18) node[midway] {\iflangja{ゲスト OS}{guest OS}};
  \draw[hv]   (34,6)  rectangle (54,12) node[midway] {\iflangja{ハイパーバイザ}{hypervisor}};
  \draw[->, thick, ln-caution] (54.3,15) to[out=0, in=0, looseness=2.2]
    node[right=0.8mm, text=ln-caution] {\iflangja{}{trap}} (54.3,9);
  % The Japanese label is wider than the gap between panels (b) and (c),
  % so it is set vertically beside the arc instead of next to it.
  \iflangja{\node[text=ln-caution, rotate=90] at (60.3,12) {トラップ};}{}
  \node[pcap] at (44,-1) {\iflangja{(b) リング 0 の\\ハイパーバイザ}{(b) hypervisor\\in ring 0}};
  % (c) root / non-root
  \node[hdr] at (74,33) {\iflangja{非ルートモード}{non-root mode}};
  \node[hdr] at (108,33) {\iflangja{ルートモード}{root mode}};
  \draw[used] (64,24) rectangle (84,30) node[midway] {\iflangja{アプリケーション}{applications}};
  \draw[free] (64,18) rectangle (84,24);
  \draw[free] (64,12) rectangle (84,18);
  \draw[used] (64,6)  rectangle (84,12) node[midway] {\iflangja{ゲスト OS}{guest OS}};
  \draw[free] (98,24) rectangle (118,30);
  \draw[free] (98,18) rectangle (118,24);
  \draw[free] (98,12) rectangle (118,18);
  \draw[hv]   (98,6)  rectangle (118,12) node[midway] {\iflangja{ハイパーバイザ}{hypervisor}};
  \draw[->, thick, ln-caution] (84.3,10.2) -- (97.7,10.2);
  \draw[->, thick, ln-tip] (97.7,7.8) -- (84.3,7.8);
  \node[text=ln-caution, anchor=south] at (91,12.2) {VM exit};
  \node[text=ln-tip, anchor=north] at (91,5.8) {VM entry};
  \node[pcap] at (91,-1) {\iflangja{(c) ルートモードと非ルートモード}{(c) root and non-root modes}};
\end{tikzpicture}
   (width ~118 mm)
\begin{tikzpicture}[x=1mm, y=1mm, font=\scriptsize\sffamily,
  >={Stealth[length=1.5mm]},
  used/.style={draw, fill=ln-accent!12},
  hv/.style={draw, fill=ln-accent!30},
  free/.style={draw=ln-rule, fill=white},
  rl/.style={font=\scriptsize\sffamily, anchor=east, text=ln-muted},
  hdr/.style={font=\scriptsize\sffamily\bfseries},
  pcap/.style={font=\footnotesize\sffamily, align=center, anchor=north}]
  % row labels
  \foreach \r/\y in {3/27, 2/21, 1/15, 0/9}
    \node[rl] at (9.5,\y) {\iflangja{リング \r}{ring \r}};
  % (a) native
  \draw[used] (11,24) rectangle (31,30) node[midway] {\iflangja{アプリケーション}{applications}};
  \draw[free] (11,18) rectangle (31,24);
  \draw[free] (11,12) rectangle (31,18);
  \draw[used] (11,6)  rectangle (31,12) node[midway] {\iflangja{OS カーネル}{OS kernel}};
  \node[pcap] at (21,-1) {\iflangja{(a) 仮想化なし}{(a) native}};
  % (b) hypervisor in ring 0
  \draw[used] (34,24) rectangle (54,30) node[midway] {\iflangja{アプリケーション}{applications}};
  \draw[free] (34,18) rectangle (54,24);
  \draw[used] (34,12) rectangle (54,18) node[midway] {\iflangja{ゲスト OS}{guest OS}};
  \draw[hv]   (34,6)  rectangle (54,12) node[midway] {\iflangja{ハイパーバイザ}{hypervisor}};
  \draw[->, thick, ln-caution] (54.3,15) to[out=0, in=0, looseness=2.2]
    node[right=0.8mm, text=ln-caution] {\iflangja{}{trap}} (54.3,9);
  % The Japanese label is wider than the gap between panels (b) and (c),
  % so it is set vertically beside the arc instead of next to it.
  \iflangja{\node[text=ln-caution, rotate=90] at (60.3,12) {トラップ};}{}
  \node[pcap] at (44,-1) {\iflangja{(b) リング 0 の\\ハイパーバイザ}{(b) hypervisor\\in ring 0}};
  % (c) root / non-root
  \node[hdr] at (74,33) {\iflangja{非ルートモード}{non-root mode}};
  \node[hdr] at (108,33) {\iflangja{ルートモード}{root mode}};
  \draw[used] (64,24) rectangle (84,30) node[midway] {\iflangja{アプリケーション}{applications}};
  \draw[free] (64,18) rectangle (84,24);
  \draw[free] (64,12) rectangle (84,18);
  \draw[used] (64,6)  rectangle (84,12) node[midway] {\iflangja{ゲスト OS}{guest OS}};
  \draw[free] (98,24) rectangle (118,30);
  \draw[free] (98,18) rectangle (118,24);
  \draw[free] (98,12) rectangle (118,18);
  \draw[hv]   (98,6)  rectangle (118,12) node[midway] {\iflangja{ハイパーバイザ}{hypervisor}};
  \draw[->, thick, ln-caution] (84.3,10.2) -- (97.7,10.2);
  \draw[->, thick, ln-tip] (97.7,7.8) -- (84.3,7.8);
  \node[text=ln-caution, anchor=south] at (91,12.2) {VM exit};
  \node[text=ln-tip, anchor=north] at (91,5.8) {VM entry};
  \node[pcap] at (91,-1) {\iflangja{(c) ルートモードと非ルートモード}{(c) root and non-root modes}};
\end{tikzpicture}

  \caption{Use of the x86 protection rings.  (a) Native execution.
    (b) Virtualization without hardware support: the guest OS is
    deprivileged to ring~1, and its privileged operations trap to the
    hypervisor.  (c) Processors with virtualization support: the guest OS
    runs in ring~0 of non-root mode, and a VM exit transfers control to
    the hypervisor in root mode.}
  \label{fig:l12-rings}
\end{figure}

Processors with virtualization support, including current x86-64
processors, go further: they provide two modes of operation, each with its
own four rings (\cref{fig:l12-rings}c).

\begin{definition}[Root mode and non-root mode]
  In \term{root mode}[ルートモード] the processor runs the hypervisor, in
  ring~0, with full control of the hardware.  In
  \term{non-root mode}[非ルートモード] it runs the VMs, with the guest OS
  in ring~0 and applications in ring~3, and privileged operations are
  restricted.
\end{definition}

The guest OS thus runs at the privilege level for which it was written.  An
attempt by the guest OS to perform a privileged operation causes a trap
called a \term{VM exit}[VM exit], which switches the processor to root mode
and passes control to the hypervisor.  When the hypervisor has handled the
operation, a \emph{VM entry}\index{VM entry} switches the processor back to
non-root mode and resumes the VM.

\section{Processor virtualization}
\label{sec:l12-cpu}

To meet the performance requirement of \cref{sec:l12-what}, the hypervisor
does not interpret the instructions of a VM but lets them execute directly
on the hardware.\source{P3L6, processor virtualization; Popek and Goldberg
\cite{popek1974}; OSTEP app.~B.}  The question is what must happen when such
an instruction is privileged.

\begin{definition}[Trap-and-emulate]
  In \term{trap-and-emulate}[トラップアンドエミュレート] virtualization,
  the instructions of a guest execute directly on the CPU, with the guest
  OS deprivileged.  Non-privileged instructions run at hardware speed.  A
  privileged instruction traps to the hypervisor, which determines what to
  do: if the operation is illegal for the VM, the hypervisor terminates the
  VM; if it is legal, the hypervisor emulates the behaviour that the guest
  OS expects from the hardware and resumes the guest.
\end{definition}

\Cref{fig:l12-trap-emulate} summarizes the decisions.  The emulation acts
on the virtual state of the VM, not on the physical hardware.  When a guest
OS disables interrupts, for example, the hypervisor records that interrupts
are disabled for this VM and holds back the interrupts destined for it,
while physical interrupts remain enabled so that the hypervisor and the
other VMs continue to receive theirs.  Every trap costs a transition into
the hypervisor, so the overhead of trap-and-emulate grows with the number
of privileged operations that the guest performs.

\begin{figure}[htbp]
  \centering
  % Trap-and-emulate: non-privileged guest instructions run directly on the
% CPU; privileged ones trap to the hypervisor, which emulates legal
% operations and terminates the VM on illegal ones.
% Usage: % Trap-and-emulate: non-privileged guest instructions run directly on the
% CPU; privileged ones trap to the hypervisor, which emulates legal
% operations and terminates the VM on illegal ones.
% Usage: % Trap-and-emulate: non-privileged guest instructions run directly on the
% CPU; privileged ones trap to the hypervisor, which emulates legal
% operations and terminates the VM on illegal ones.
% Usage: \input{figures/l12-trap-emulate}   (width ~118 mm)
\begin{tikzpicture}[x=1mm, y=1mm, font=\scriptsize\sffamily,
  >={Stealth[length=1.5mm]},
  stp/.style={draw, rounded corners=1.5pt, fill=white, align=center,
               minimum height=9mm, inner sep=1pt},
  hyp/.style={stp, fill=ln-accent!15},
  dec/.style={draw, diamond, aspect=2, fill=ln-box, inner sep=0.3pt,
              align=center},
  yn/.style={font=\scriptsize\sffamily\itshape, text=ln-muted}]
  \node[stp, minimum width=24mm] (ins) at (12,28)
    {\iflangja{ゲストが命令を\\実行する}{guest executes\\an instruction}};
  \node[dec] (priv) at (42,28) {\iflangja{特権命令?}{privileged?}};
  \node[stp, minimum width=38mm] (hw) at (93,28)
    {\iflangja{CPU 上で直接実行\\（ハードウェアの速度）}{runs directly on the CPU\\(hardware speed)}};
  \node[hyp, minimum width=24mm] (trap) at (42,13)
    {\iflangja{ハイパーバイザへ\\トラップ}{trap to the\\hypervisor}};
  \node[dec] (legal) at (70,13) {\iflangja{正当?}{legal?}};
  \node[hyp, minimum width=30mm] (emu) at (102,13)
    {\iflangja{期待された効果を仮想\\状態上でエミュレート}{emulate the expected\\effect on virtual state}};
  \node[hyp, minimum width=26mm, minimum height=7mm] (term) at (70,-2)
    {\iflangja{VM を終了させる}{terminate the VM}};
  \draw[->] (ins) -- (priv);
  \draw[->] (priv) -- node[yn, above] {\iflangja{いいえ}{no}} (hw);
  \draw[->] (priv) -- node[yn, right] {\iflangja{はい}{yes}} (trap);
  \draw[->] (trap) -- (legal);
  \draw[->] (legal) -- node[yn, above] {\iflangja{はい}{yes}} (emu);
  \draw[->] (legal) -- node[yn, right] {\iflangja{いいえ}{no}} (term);
  \draw[->, densely dashed, ln-muted] (emu.east) -- ++(2.5,0) |- (12,38)
    -- (ins.north);
  \draw[densely dashed, ln-muted] (hw.north) -- (93,38);
  \node[yn, anchor=south] at (52,38) {\iflangja{次の命令（ゲストを再開）}{next instruction (guest resumes)}};
\end{tikzpicture}
   (width ~118 mm)
\begin{tikzpicture}[x=1mm, y=1mm, font=\scriptsize\sffamily,
  >={Stealth[length=1.5mm]},
  stp/.style={draw, rounded corners=1.5pt, fill=white, align=center,
               minimum height=9mm, inner sep=1pt},
  hyp/.style={stp, fill=ln-accent!15},
  dec/.style={draw, diamond, aspect=2, fill=ln-box, inner sep=0.3pt,
              align=center},
  yn/.style={font=\scriptsize\sffamily\itshape, text=ln-muted}]
  \node[stp, minimum width=24mm] (ins) at (12,28)
    {\iflangja{ゲストが命令を\\実行する}{guest executes\\an instruction}};
  \node[dec] (priv) at (42,28) {\iflangja{特権命令?}{privileged?}};
  \node[stp, minimum width=38mm] (hw) at (93,28)
    {\iflangja{CPU 上で直接実行\\（ハードウェアの速度）}{runs directly on the CPU\\(hardware speed)}};
  \node[hyp, minimum width=24mm] (trap) at (42,13)
    {\iflangja{ハイパーバイザへ\\トラップ}{trap to the\\hypervisor}};
  \node[dec] (legal) at (70,13) {\iflangja{正当?}{legal?}};
  \node[hyp, minimum width=30mm] (emu) at (102,13)
    {\iflangja{期待された効果を仮想\\状態上でエミュレート}{emulate the expected\\effect on virtual state}};
  \node[hyp, minimum width=26mm, minimum height=7mm] (term) at (70,-2)
    {\iflangja{VM を終了させる}{terminate the VM}};
  \draw[->] (ins) -- (priv);
  \draw[->] (priv) -- node[yn, above] {\iflangja{いいえ}{no}} (hw);
  \draw[->] (priv) -- node[yn, right] {\iflangja{はい}{yes}} (trap);
  \draw[->] (trap) -- (legal);
  \draw[->] (legal) -- node[yn, above] {\iflangja{はい}{yes}} (emu);
  \draw[->] (legal) -- node[yn, right] {\iflangja{いいえ}{no}} (term);
  \draw[->, densely dashed, ln-muted] (emu.east) -- ++(2.5,0) |- (12,38)
    -- (ins.north);
  \draw[densely dashed, ln-muted] (hw.north) -- (93,38);
  \node[yn, anchor=south] at (52,38) {\iflangja{次の命令（ゲストを再開）}{next instruction (guest resumes)}};
\end{tikzpicture}
   (width ~118 mm)
\begin{tikzpicture}[x=1mm, y=1mm, font=\scriptsize\sffamily,
  >={Stealth[length=1.5mm]},
  stp/.style={draw, rounded corners=1.5pt, fill=white, align=center,
               minimum height=9mm, inner sep=1pt},
  hyp/.style={stp, fill=ln-accent!15},
  dec/.style={draw, diamond, aspect=2, fill=ln-box, inner sep=0.3pt,
              align=center},
  yn/.style={font=\scriptsize\sffamily\itshape, text=ln-muted}]
  \node[stp, minimum width=24mm] (ins) at (12,28)
    {\iflangja{ゲストが命令を\\実行する}{guest executes\\an instruction}};
  \node[dec] (priv) at (42,28) {\iflangja{特権命令?}{privileged?}};
  \node[stp, minimum width=38mm] (hw) at (93,28)
    {\iflangja{CPU 上で直接実行\\（ハードウェアの速度）}{runs directly on the CPU\\(hardware speed)}};
  \node[hyp, minimum width=24mm] (trap) at (42,13)
    {\iflangja{ハイパーバイザへ\\トラップ}{trap to the\\hypervisor}};
  \node[dec] (legal) at (70,13) {\iflangja{正当?}{legal?}};
  \node[hyp, minimum width=30mm] (emu) at (102,13)
    {\iflangja{期待された効果を仮想\\状態上でエミュレート}{emulate the expected\\effect on virtual state}};
  \node[hyp, minimum width=26mm, minimum height=7mm] (term) at (70,-2)
    {\iflangja{VM を終了させる}{terminate the VM}};
  \draw[->] (ins) -- (priv);
  \draw[->] (priv) -- node[yn, above] {\iflangja{いいえ}{no}} (hw);
  \draw[->] (priv) -- node[yn, right] {\iflangja{はい}{yes}} (trap);
  \draw[->] (trap) -- (legal);
  \draw[->] (legal) -- node[yn, above] {\iflangja{はい}{yes}} (emu);
  \draw[->] (legal) -- node[yn, right] {\iflangja{いいえ}{no}} (term);
  \draw[->, densely dashed, ln-muted] (emu.east) -- ++(2.5,0) |- (12,38)
    -- (ins.north);
  \draw[densely dashed, ln-muted] (hw.north) -- (93,38);
  \node[yn, anchor=south] at (52,38) {\iflangja{次の命令（ゲストを再開）}{next instruction (guest resumes)}};
\end{tikzpicture}

  \caption{Trap-and-emulate.  Only privileged instructions involve the
    hypervisor; legal ones are emulated on the virtual state of the VM.}
  \label{fig:l12-trap-emulate}
\end{figure}

\begin{note}
  Trap-and-emulate is correct only if every instruction whose effect the
  hypervisor must control actually traps.  Popek and Goldberg
  \cite{popek1974} call such instructions \emph{sensitive}: they change or
  reveal the configuration of the machine, or behave differently depending
  on the privilege level.  An architecture can be virtualized in this way
  if every sensitive instruction is privileged, that is, traps when
  executed outside the most privileged level.
\end{note}

\section{x86 virtualization before hardware support}
\label{sec:l12-x86}

Before 2005, x86 processors had four rings but no root and non-root modes,
so a hypervisor ran in ring~0 and deprivileged the guest OS to
ring~1.\source{P3L6, x86 virtualization in the past; Robin and Irvine
\cite{robin2000}.}  The architecture did not satisfy the condition of
trap-and-emulate: Robin and Irvine \cite{robin2000} identified 17 sensitive
instructions that do not trap when executed in a less privileged ring.  Executed in
ring~1, they fail silently or reveal machine state that the hypervisor
intends to hide, and trap-and-emulate alone cannot virtualize the
processor correctly.

The best-known examples are \texttt{PUSHF} and \texttt{POPF}, which save
the flags register on the stack and restore it.  The flags register
contains the interrupt-enable flag IF, and kernels commonly re-enable
interrupts at the end of a section that ran with interrupts disabled by
restoring the saved flags with \texttt{POPF}.  In ring~0, \texttt{POPF}
restores IF.  In ring~1 the processor ignores the IF bit of the restored
value, without a trap.\margin{Precisely, \texttt{POPF} leaves IF unchanged
whenever the current ring is numerically greater than the I/O privilege
level in the flags register, which the hypervisor sets to~0.  \texttt{CLI}
traps under the same condition.}  The hypervisor is not invoked, so it does
not learn that the guest wanted to change the interrupt state and cannot
emulate the change; the guest OS receives no error and assumes that the
change succeeded.  \Cref{tab:l12-popf} traces the consequence.

\begin{table}[htbp]
  \centering
  \caption{A guest kernel in ring~1 saves the flags, disables interrupts
    and later restores the flags.  \texttt{CLI} traps and is emulated;
    \texttt{POPF} does not trap.  \emph{Virtual IF} is the interrupt flag
    that the hypervisor records for the VM; physical interrupts stay
    enabled throughout.}
  \label{tab:l12-popf}
  \small
  \begin{tabular}{@{}clccc@{}}
    \toprule
    Step & Guest kernel executes & Traps & Virtual IF & Guest assumes \\
    \midrule
    1 & \texttt{PUSHF} (saves IF $=1$)     & no  & 1 & enabled \\
    2 & \texttt{CLI}                         & yes & 0 & disabled \\
    3 & critical section                     & no  & 0 & disabled \\
    4 & \texttt{POPF} (restores IF $=1$)   & no  & 0 & enabled \\
    5 & waits for a device interrupt         & no  & 0 & enabled \\
    \bottomrule
  \end{tabular}
\end{table}

From step~4 on, the hypervisor considers interrupts disabled for the VM and
withholds its interrupts, while the guest OS waits for an interrupt that it
believes to be enabled: the guest hangs.  Two remedies were developed in
software, one that leaves the guest OS unmodified and one that modifies it.

\section{Binary translation}
\label{sec:l12-bt}

The first remedy rewrites the binary code of the guest so that it never
executes the problematic instructions.\source{P3L6, binary translation;
Bugnion et al.\ \cite{bugnion2012}; Silberschatz et al.\ ch.~18.}  The
approach was pioneered by Mendel Rosenblum's group at Stanford University
and commercialized by VMware \cite{bugnion2012}.  Its goal is to run the
guest OS without any modification.

\begin{definition}[Full virtualization]
  Virtualization in which the guest operating system runs unmodified, as
  it would on the physical machine, is called
  \term{full virtualization}[完全仮想化].
\end{definition}

\begin{definition}[Binary translation]
  In dynamic \term{binary translation}[バイナリ変換] the hypervisor rewrites
  guest code at run time, just before it executes.  Code that contains
  problematic instructions is translated into an alternative instruction
  sequence with the intended effect; all other code runs unchanged at
  hardware speed.
\end{definition}

The translation proceeds block by block (\cref{fig:l12-bt}).
\begin{enumerate}
  \item Inspect the block of code that is about to be executed, for
    example the instructions up to the next branch.
  \item If the block contains problematic instructions, translate them into
    a sequence that emulates the desired behaviour, for example a call to
    an emulation routine of the hypervisor.  The translation may even avoid
    a trap where the original instruction would have trapped.
  \item Otherwise execute the block as it is, at hardware speed.
\end{enumerate}
Translated blocks are kept in a \emph{translation cache}, so the cost of
translating a block is paid once and amortized over its later executions.
The translation must be dynamic because which code executes, and where
indirect branches lead, is known only at run time.%
\margin{In VMware's design only guest kernel code is translated;
applications in ring~3 run directly \cite{bugnion2012}.}

\begin{figure}[htbp]
  \centering
  % Dynamic binary translation with a translation cache: a guest kernel code
% block is inspected and translated the first time it is about to run;
% later executions reuse the cached translation.
% Usage: % Dynamic binary translation with a translation cache: a guest kernel code
% block is inspected and translated the first time it is about to run;
% later executions reuse the cached translation.
% Usage: % Dynamic binary translation with a translation cache: a guest kernel code
% block is inspected and translated the first time it is about to run;
% later executions reuse the cached translation.
% Usage: \input{figures/l12-binary-translation}   (width ~119 mm)
\begin{tikzpicture}[x=1mm, y=1mm, font=\scriptsize\sffamily,
  >={Stealth[length=1.5mm]},
  stp/.style={draw, rounded corners=1.5pt, fill=white, align=center,
               minimum height=9mm, inner sep=1pt},
  hyp/.style={stp, fill=ln-accent!15},
  dec/.style={draw, diamond, aspect=2.2, fill=ln-box, inner sep=0.3pt,
              align=center},
  yn/.style={font=\scriptsize\sffamily\itshape, text=ln-muted}]
  \node[stp, minimum width=22mm] (blk) at (11,26)
    {\iflangja{次のゲスト\\カーネルコード\\ブロック}{next guest\\kernel code\\block}};
  \node[dec] (hit) at (42,26)
    {\iflangja{変換キャッシュ\\にある?}{in translation\\cache?}};
  \node[stp, minimum width=32mm] (exe) at (103,26)
    {\iflangja{変換済みブロックを\\ハードウェアの速度で実行}{execute translated block\\at hardware speed}};
  \node[hyp, minimum width=22mm] (insp) at (42,8)
    {\iflangja{命令を\\検査する}{inspect the\\instructions}};
  \node[hyp, minimum width=26mm] (tr) at (72,8)
    {\iflangja{問題のある命令を\\別の命令列に置換}{replace problematic\\instructions}};
  \node[hyp, minimum width=28mm] (st) at (103,8)
    {\iflangja{変換キャッシュに\\格納する}{store the block in\\the translation cache}};
  \draw[->] (blk) -- (hit);
  \draw[->] (hit) -- node[yn, above] {\iflangja{はい}{yes}} (exe);
  \draw[->] (hit) -- node[yn, right] {\iflangja{いいえ}{no}} (insp);
  \draw[->] (insp) -- (tr);
  \draw[->] (tr) -- (st);
  \draw[->] (st) -- (exe);
  \draw[->, densely dashed, ln-muted] (exe.north) -- (103,39) -- (11,39)
    -- (blk.north);
  \node[yn, anchor=south] at (57,39) {\iflangja{次のブロックへ}{continue with the next block}};
\end{tikzpicture}
   (width ~119 mm)
\begin{tikzpicture}[x=1mm, y=1mm, font=\scriptsize\sffamily,
  >={Stealth[length=1.5mm]},
  stp/.style={draw, rounded corners=1.5pt, fill=white, align=center,
               minimum height=9mm, inner sep=1pt},
  hyp/.style={stp, fill=ln-accent!15},
  dec/.style={draw, diamond, aspect=2.2, fill=ln-box, inner sep=0.3pt,
              align=center},
  yn/.style={font=\scriptsize\sffamily\itshape, text=ln-muted}]
  \node[stp, minimum width=22mm] (blk) at (11,26)
    {\iflangja{次のゲスト\\カーネルコード\\ブロック}{next guest\\kernel code\\block}};
  \node[dec] (hit) at (42,26)
    {\iflangja{変換キャッシュ\\にある?}{in translation\\cache?}};
  \node[stp, minimum width=32mm] (exe) at (103,26)
    {\iflangja{変換済みブロックを\\ハードウェアの速度で実行}{execute translated block\\at hardware speed}};
  \node[hyp, minimum width=22mm] (insp) at (42,8)
    {\iflangja{命令を\\検査する}{inspect the\\instructions}};
  \node[hyp, minimum width=26mm] (tr) at (72,8)
    {\iflangja{問題のある命令を\\別の命令列に置換}{replace problematic\\instructions}};
  \node[hyp, minimum width=28mm] (st) at (103,8)
    {\iflangja{変換キャッシュに\\格納する}{store the block in\\the translation cache}};
  \draw[->] (blk) -- (hit);
  \draw[->] (hit) -- node[yn, above] {\iflangja{はい}{yes}} (exe);
  \draw[->] (hit) -- node[yn, right] {\iflangja{いいえ}{no}} (insp);
  \draw[->] (insp) -- (tr);
  \draw[->] (tr) -- (st);
  \draw[->] (st) -- (exe);
  \draw[->, densely dashed, ln-muted] (exe.north) -- (103,39) -- (11,39)
    -- (blk.north);
  \node[yn, anchor=south] at (57,39) {\iflangja{次のブロックへ}{continue with the next block}};
\end{tikzpicture}
   (width ~119 mm)
\begin{tikzpicture}[x=1mm, y=1mm, font=\scriptsize\sffamily,
  >={Stealth[length=1.5mm]},
  stp/.style={draw, rounded corners=1.5pt, fill=white, align=center,
               minimum height=9mm, inner sep=1pt},
  hyp/.style={stp, fill=ln-accent!15},
  dec/.style={draw, diamond, aspect=2.2, fill=ln-box, inner sep=0.3pt,
              align=center},
  yn/.style={font=\scriptsize\sffamily\itshape, text=ln-muted}]
  \node[stp, minimum width=22mm] (blk) at (11,26)
    {\iflangja{次のゲスト\\カーネルコード\\ブロック}{next guest\\kernel code\\block}};
  \node[dec] (hit) at (42,26)
    {\iflangja{変換キャッシュ\\にある?}{in translation\\cache?}};
  \node[stp, minimum width=32mm] (exe) at (103,26)
    {\iflangja{変換済みブロックを\\ハードウェアの速度で実行}{execute translated block\\at hardware speed}};
  \node[hyp, minimum width=22mm] (insp) at (42,8)
    {\iflangja{命令を\\検査する}{inspect the\\instructions}};
  \node[hyp, minimum width=26mm] (tr) at (72,8)
    {\iflangja{問題のある命令を\\別の命令列に置換}{replace problematic\\instructions}};
  \node[hyp, minimum width=28mm] (st) at (103,8)
    {\iflangja{変換キャッシュに\\格納する}{store the block in\\the translation cache}};
  \draw[->] (blk) -- (hit);
  \draw[->] (hit) -- node[yn, above] {\iflangja{はい}{yes}} (exe);
  \draw[->] (hit) -- node[yn, right] {\iflangja{いいえ}{no}} (insp);
  \draw[->] (insp) -- (tr);
  \draw[->] (tr) -- (st);
  \draw[->] (st) -- (exe);
  \draw[->, densely dashed, ln-muted] (exe.north) -- (103,39) -- (11,39)
    -- (blk.north);
  \node[yn, anchor=south] at (57,39) {\iflangja{次のブロックへ}{continue with the next block}};
\end{tikzpicture}

  \caption{Dynamic binary translation with a translation cache.  A block is
    inspected and translated only the first time it is about to run.}
  \label{fig:l12-bt}
\end{figure}

\needspace{8\baselineskip}
\begin{example}
  A block of guest kernel code executes natively in \SI{5}{\micro\second}.
  Translating it takes \SI{2}{\milli\second}, and because one instruction
  is replaced by a call into the hypervisor, the translated block takes
  \SI{6}{\micro\second}.  If the block runs \num{20000} times, translating
  it once and caching the result costs
  $\SI{2}{\milli\second} + \num{20000}\times\SI{6}{\micro\second} =
  \SI{122}{\milli\second}$, against \SI{100}{\milli\second} natively, an
  overhead of \SI{22}{\percent}.  Translating it anew before every
  execution would cost
  $\num{20000}\times(\SI{2}{\milli\second}+\SI{6}{\micro\second}) \approx
  \SI{40.1}{\second}$, about 400 times the native time.
\end{example}

\section{Paravirtualization}
\label{sec:l12-pv}

The second remedy gives up on unmodified guests in exchange for
performance.\source{P3L6, paravirtualization; Barham et al.\
\cite{barham2003}.}  The guest OS is ported to the hypervisor, much as an
operating system is ported to a new hardware platform.

\begin{definition}[Paravirtualization]
  In \term{paravirtualization}[準仮想化] the guest operating system is
  modified so that it knows it runs in a VM.  Instead of executing
  privileged operations and relying on traps, it requests them from the
  hypervisor explicitly.
\end{definition}

An explicit request of a guest OS to the hypervisor is a
\term{hypercall}[ハイパーコール].  It plays the role for a guest OS that a
system call plays for an application (Lesson~1): the guest places the
context information---the arguments---in a well-defined location,
specifies which hypercall it requests, and traps into the hypervisor, which
performs the operation and returns control to the guest.  The hypervisor
therefore does not have to infer the intended operation from individual
trapped instructions, and one hypercall can carry several operations, as
\cref{sec:l12-memory} shows.  Only the guest OS is modified; applications
run unchanged, because the interface that the guest OS offers them stays
the same \cite{barham2003}.

Xen, developed at the University of Cambridge and later commercialized by
XenSource and Citrix, is the best-known paravirtualizing hypervisor.
\Cref{tab:l12-cpu} in \cref{sec:l12-hw} compares paravirtualization with
the other approaches to processor virtualization.

\begin{keypoint}
  Trap-and-emulate requires every sensitive instruction to trap, which
  x86 processors before 2005 did not guarantee.  Binary translation
  rewrites the problematic instructions at run time and keeps the guest OS
  unmodified; paravirtualization modifies the guest OS to issue hypercalls
  and gains performance.
\end{keypoint}

\section{Memory virtualization}
\label{sec:l12-memory}

\subsection{Full virtualization}

A guest OS expects what an operating system finds on a physical machine:
physical memory that is contiguous and starts at address~0, which it
manages with page tables as described in Lesson~8.\source{P3L6, memory
virtualization: full, paravirtualized; OSTEP app.~B; Barham et al.\
\cite{barham2003}.}  The hypervisor must give every VM this illusion while
dividing the memory of the machine among the VMs, and it should still let
the hardware MMU and TLB translate addresses.  Three kinds of addresses are
therefore distinguished:
\begin{description}
  \item[Virtual addresses] are used by the software in a VM, as in
    Lesson~2.
  \item[Physical addresses] are the addresses that the guest OS believes to
    be those of physical memory.
  \item[Machine addresses] are the actual addresses of the memory of the
    machine.  A \term{machine address}[マシンアドレス] is known only to the
    hypervisor under full virtualization.
\end{description}
Page numbers make the same distinction between virtual page numbers,
physical frame numbers and machine frame numbers.  The guest OS maintains
page tables that map virtual to physical addresses, and the hypervisor
maintains a mapping from physical to machine addresses
(\cref{fig:l12-memory}).

\begin{figure}[htbp]
  \centering
  % Memory virtualization: virtual, physical and machine addresses.  The
% guest page table maps VA to PA, the hypervisor maps PA to MA; the shadow
% page table maps VA directly to MA and is the table the MMU uses.
% Usage: % Memory virtualization: virtual, physical and machine addresses.  The
% guest page table maps VA to PA, the hypervisor maps PA to MA; the shadow
% page table maps VA directly to MA and is the table the MMU uses.
% Usage: % Memory virtualization: virtual, physical and machine addresses.  The
% guest page table maps VA to PA, the hypervisor maps PA to MA; the shadow
% page table maps VA directly to MA and is the table the MMU uses.
% Usage: \input{figures/l12-memory}   (width ~120 mm)
\begin{tikzpicture}[x=1mm, y=1mm, font=\scriptsize\sffamily,
  >={Stealth[length=1.5mm]},
  addr/.style={draw, rounded corners=4pt, fill=white, align=center,
               minimum height=9mm, inner sep=1pt},
  tab/.style={draw, fill=ln-accent!12, align=center, minimum height=10mm,
              minimum width=22mm, inner sep=1pt},
  htab/.style={tab, fill=ln-accent!30},
  own/.style={font=\scriptsize\sffamily\bfseries, text=ln-accent},
  lab/.style={font=\scriptsize\sffamily, text=ln-muted, align=center}]
  \node[addr, minimum width=17mm] (va) at (8.5,22)
    {\iflangja{仮想\\アドレス}{virtual\\address}};
  \node[tab] (gpt) at (34,22)
    {\iflangja{ゲストの\\ページテーブル}{guest\\page table}};
  \node[addr, minimum width=20mm] (pa) at (60,22)
    {\iflangja{物理\\アドレス}{physical\\address}};
  \node[htab] (p2m) at (86,22)
    {\iflangja{ハイパーバイザの\\PA$\to$MA 対応表}{hypervisor's\\PA$\to$MA map}};
  \node[addr, minimum width=17mm] (ma) at (111.5,22)
    {\iflangja{マシン\\アドレス}{machine\\address}};
  \draw[->] (va) -- (gpt);
  \draw[->] (gpt) -- (pa);
  \draw[->] (pa) -- (p2m);
  \draw[->] (p2m) -- (ma);
  \node[own, anchor=south] at (gpt.north) {\iflangja{ゲスト OS が管理}{kept by the guest OS}};
  \node[own, anchor=south] at (p2m.north) {\iflangja{ハイパーバイザが管理}{kept by the hypervisor}};
  \draw[ln-muted] (23,31.5) -- (23,33.5) -- (97,33.5) -- (97,31.5);
  \node[lab, anchor=south] at (60,33.7) {\iflangja{2 回の変換}{two translations per address}};
  % shadow page table
  \node[htab, minimum width=40mm, minimum height=8mm] (spt) at (60,4)
    {\iflangja{シャドウページテーブル（VA$\to$MA）}{shadow page table (VA$\to$MA)}};
  \draw[->, very thick, ln-accent] (va.south) |- (spt.west);
  \draw[->, very thick, ln-accent] (spt.east) -| (ma.south);
  \draw[->, densely dashed, ln-muted] (gpt.south) -- (44,8.3);
  \draw[->, densely dashed, ln-muted] (p2m.south) -- (76,8.3);
  \node[lab] at (60,12.8) {\iflangja{ハイパーバイザが一貫性を保つ}{kept consistent by the hypervisor}};
  \node[lab, anchor=north] at (60,-0.4)
    {\iflangja{1 回の変換．MMU と TLB が使用}{one translation, used by the MMU and TLB}};
\end{tikzpicture}
   (width ~120 mm)
\begin{tikzpicture}[x=1mm, y=1mm, font=\scriptsize\sffamily,
  >={Stealth[length=1.5mm]},
  addr/.style={draw, rounded corners=4pt, fill=white, align=center,
               minimum height=9mm, inner sep=1pt},
  tab/.style={draw, fill=ln-accent!12, align=center, minimum height=10mm,
              minimum width=22mm, inner sep=1pt},
  htab/.style={tab, fill=ln-accent!30},
  own/.style={font=\scriptsize\sffamily\bfseries, text=ln-accent},
  lab/.style={font=\scriptsize\sffamily, text=ln-muted, align=center}]
  \node[addr, minimum width=17mm] (va) at (8.5,22)
    {\iflangja{仮想\\アドレス}{virtual\\address}};
  \node[tab] (gpt) at (34,22)
    {\iflangja{ゲストの\\ページテーブル}{guest\\page table}};
  \node[addr, minimum width=20mm] (pa) at (60,22)
    {\iflangja{物理\\アドレス}{physical\\address}};
  \node[htab] (p2m) at (86,22)
    {\iflangja{ハイパーバイザの\\PA$\to$MA 対応表}{hypervisor's\\PA$\to$MA map}};
  \node[addr, minimum width=17mm] (ma) at (111.5,22)
    {\iflangja{マシン\\アドレス}{machine\\address}};
  \draw[->] (va) -- (gpt);
  \draw[->] (gpt) -- (pa);
  \draw[->] (pa) -- (p2m);
  \draw[->] (p2m) -- (ma);
  \node[own, anchor=south] at (gpt.north) {\iflangja{ゲスト OS が管理}{kept by the guest OS}};
  \node[own, anchor=south] at (p2m.north) {\iflangja{ハイパーバイザが管理}{kept by the hypervisor}};
  \draw[ln-muted] (23,31.5) -- (23,33.5) -- (97,33.5) -- (97,31.5);
  \node[lab, anchor=south] at (60,33.7) {\iflangja{2 回の変換}{two translations per address}};
  % shadow page table
  \node[htab, minimum width=40mm, minimum height=8mm] (spt) at (60,4)
    {\iflangja{シャドウページテーブル（VA$\to$MA）}{shadow page table (VA$\to$MA)}};
  \draw[->, very thick, ln-accent] (va.south) |- (spt.west);
  \draw[->, very thick, ln-accent] (spt.east) -| (ma.south);
  \draw[->, densely dashed, ln-muted] (gpt.south) -- (44,8.3);
  \draw[->, densely dashed, ln-muted] (p2m.south) -- (76,8.3);
  \node[lab] at (60,12.8) {\iflangja{ハイパーバイザが一貫性を保つ}{kept consistent by the hypervisor}};
  \node[lab, anchor=north] at (60,-0.4)
    {\iflangja{1 回の変換．MMU と TLB が使用}{one translation, used by the MMU and TLB}};
\end{tikzpicture}
   (width ~120 mm)
\begin{tikzpicture}[x=1mm, y=1mm, font=\scriptsize\sffamily,
  >={Stealth[length=1.5mm]},
  addr/.style={draw, rounded corners=4pt, fill=white, align=center,
               minimum height=9mm, inner sep=1pt},
  tab/.style={draw, fill=ln-accent!12, align=center, minimum height=10mm,
              minimum width=22mm, inner sep=1pt},
  htab/.style={tab, fill=ln-accent!30},
  own/.style={font=\scriptsize\sffamily\bfseries, text=ln-accent},
  lab/.style={font=\scriptsize\sffamily, text=ln-muted, align=center}]
  \node[addr, minimum width=17mm] (va) at (8.5,22)
    {\iflangja{仮想\\アドレス}{virtual\\address}};
  \node[tab] (gpt) at (34,22)
    {\iflangja{ゲストの\\ページテーブル}{guest\\page table}};
  \node[addr, minimum width=20mm] (pa) at (60,22)
    {\iflangja{物理\\アドレス}{physical\\address}};
  \node[htab] (p2m) at (86,22)
    {\iflangja{ハイパーバイザの\\PA$\to$MA 対応表}{hypervisor's\\PA$\to$MA map}};
  \node[addr, minimum width=17mm] (ma) at (111.5,22)
    {\iflangja{マシン\\アドレス}{machine\\address}};
  \draw[->] (va) -- (gpt);
  \draw[->] (gpt) -- (pa);
  \draw[->] (pa) -- (p2m);
  \draw[->] (p2m) -- (ma);
  \node[own, anchor=south] at (gpt.north) {\iflangja{ゲスト OS が管理}{kept by the guest OS}};
  \node[own, anchor=south] at (p2m.north) {\iflangja{ハイパーバイザが管理}{kept by the hypervisor}};
  \draw[ln-muted] (23,31.5) -- (23,33.5) -- (97,33.5) -- (97,31.5);
  \node[lab, anchor=south] at (60,33.7) {\iflangja{2 回の変換}{two translations per address}};
  % shadow page table
  \node[htab, minimum width=40mm, minimum height=8mm] (spt) at (60,4)
    {\iflangja{シャドウページテーブル（VA$\to$MA）}{shadow page table (VA$\to$MA)}};
  \draw[->, very thick, ln-accent] (va.south) |- (spt.west);
  \draw[->, very thick, ln-accent] (spt.east) -| (ma.south);
  \draw[->, densely dashed, ln-muted] (gpt.south) -- (44,8.3);
  \draw[->, densely dashed, ln-muted] (p2m.south) -- (76,8.3);
  \node[lab] at (60,12.8) {\iflangja{ハイパーバイザが一貫性を保つ}{kept consistent by the hypervisor}};
  \node[lab, anchor=north] at (60,-0.4)
    {\iflangja{1 回の変換．MMU と TLB が使用}{one translation, used by the MMU and TLB}};
\end{tikzpicture}

  \caption{Memory virtualization.  Translating through the guest page table
    and then through the hypervisor's mapping takes two steps; the shadow
    page table combines them into one mapping that the MMU uses directly.
    Hardware with extended page tables performs the two-step translation
    itself (\cref{sec:l12-hw}).}
  \label{fig:l12-memory}
\end{figure}

A first option translates every address in these two steps.  The MMU,
however, walks only one page table, and the TLB caches one translation per
address.  Performing the second translation in software on every memory
reference is far too expensive.

\begin{definition}[Shadow page table]
  A \term{shadow page table}[シャドウページテーブル] is a page table
  maintained by the hypervisor that maps the virtual addresses of a guest
  process directly to machine addresses.  The hypervisor installs it in the
  MMU in place of the guest's page table, so that the hardware translates
  each address in one step and the TLB caches the result.
\end{definition}

The guest OS continues to create and change its own page tables, which the
hardware no longer uses, and the hypervisor must keep the shadow page
tables consistent with them.
\begin{itemize}
  \item It write-protects the memory that holds the guest page tables.
    When the guest OS adds or changes a mapping, the write traps; the
    hypervisor performs the write and updates the shadow page table with
    the corresponding machine address.
  \item When the guest OS switches between its processes, it loads another
    page table into the page table base register, a privileged operation
    that traps.  The hypervisor then installs the shadow page table of the
    new process, or invalidates the current shadow entries and rebuilds
    them as they are used.
\end{itemize}
These traps make page table updates and context switches in a guest far
more costly than on a physical machine.

\begin{example}
  A guest process uses virtual pages 2 and~3.  The guest page table maps
  them to physical frames 5 and~9, which the hypervisor has backed with
  machine frames 41 and~17, so the shadow page table maps page~2 to
  machine frame~41 and page~3 to machine frame~17.  The guest OS now maps
  virtual page~4 to physical frame~12.  Its write to the protected page
  table traps; the hypervisor finds that frame~12 is backed by machine
  frame~26, performs the write, adds the entry $4 \to 26$ to the shadow
  page table and resumes the guest.  \Cref{tab:l12-shadow} shows the
  resulting mappings.
\end{example}

\begin{table}[htbp]
  \centering
  \caption{Mappings of the guest process of the example.  The last shadow
    entry is added by the hypervisor when the guest's page table update
    traps.}
  \label{tab:l12-shadow}
  \small
  \begin{tabular}{@{}cccc@{}}
    \toprule
    Virtual & Guest page table: & Hypervisor: & Shadow page table: \\
    page & physical frame & machine frame & machine frame \\
    \midrule
    2 & 5  & 41 & 41 \\
    3 & 9  & 17 & 17 \\
    4 & 12 & 26 & 26 \\
    \bottomrule
  \end{tabular}
\end{table}

\subsection{Paravirtualization}

With paravirtualization the guest OS knows that it runs in a VM, and memory
virtualization becomes a cooperation between guest and hypervisor.  The
strict requirement that physical memory be contiguous and start at
address~0 disappears.  In Xen the guest OS knows which machine frames it
has been given, and its page tables map virtual addresses directly to
machine addresses \cite{barham2003}.  The guest registers its page tables
with the hypervisor explicitly, and from then on it has only read access to
them.  To change a page table, the guest issues a hypercall, and the
hypervisor validates the change---for example, that the new mapping refers
to a machine frame that belongs to this VM---before applying it.  Because
every hypercall costs a transition into the hypervisor, the guest
\emph{batches} page table updates so that many of them share one
hypercall.  Other optimizations follow the same idea of reducing
transitions: Xen, for instance, occupies the top \SI{64}{\mega\byte} of
every guest address space, so that entering the hypervisor requires no
change of page table and no TLB flush \cite{barham2003}.

\needspace{16\baselineskip}
\begin{example}
  A paravirtualized Linux guest on Xen changes three page table entries
  with a single hypercall.  Each request gives the machine address of a
  page table entry and the new contents of that entry:
  \begin{lstlisting}[language=C, numbers=none]
struct mmu_update req[3];
int done;

for (int i = 0; i < 3; i++) {
  /* machine address of the entry to change */
  req[i].ptr = pte_maddr[i] | MMU_NORMAL_PT_UPDATE;
  /* new entry: machine frame number and flags */
  req[i].val = (mfn[i] << PAGE_SHIFT) | flags;
}
/* one transition into Xen for all three updates */
HYPERVISOR_mmu_update(req, 3, &done, DOMID_SELF);
  \end{lstlisting}
  Xen checks every request before applying it and reports in
  \texttt{done} how many updates succeeded.  Issued one by one, the same
  updates would enter the hypervisor three times.
\end{example}

On newer platforms, the hardware support of \cref{sec:l12-hw} eliminates
or greatly reduces the overheads of memory virtualization.

\section{Device virtualization}
\label{sec:l12-devices}

CPUs and memory show little diversity: the instruction set architecture
standardizes their interface, so a few techniques suffice for all
processors of a family.\source{P3L6, device virtualization; passthrough,
hypervisor-direct and split-device driver models; Silberschatz et al.\
ch.~18.}  Devices are highly diverse, and there is no standard
specification of their interface and behaviour; as Lesson~11 showed, each
type of device needs its own driver.  Three models for virtualizing devices
have emerged.  They differ in where the driver that controls the physical
device resides and in how the device accesses of a VM reach it
(\cref{fig:l12-devices}).

\begin{figure}[htbp]
  \centering
  % Device virtualization models: (a) passthrough, (b) hypervisor-direct,
% (c) split-device driver with front-end and back-end drivers.
% Usage: % Device virtualization models: (a) passthrough, (b) hypervisor-direct,
% (c) split-device driver with front-end and back-end drivers.
% Usage: % Device virtualization models: (a) passthrough, (b) hypervisor-direct,
% (c) split-device driver with front-end and back-end drivers.
% Usage: \input{figures/l12-devices}   (width ~118 mm)
\begin{tikzpicture}[x=1mm, y=1mm, font=\scriptsize\sffamily,
  >={Stealth[length=1.5mm]},
  dev/.style={draw, fill=white},
  band/.style={draw, fill=ln-accent!22},
  vm/.style={draw=ln-rule, rounded corners=2pt, fill=ln-box},
  svm/.style={draw=ln-accent, rounded corners=2pt, fill=ln-accent!8},
  gos/.style={draw, fill=ln-accent!12},
  box/.style={draw, fill=white, align=center, inner sep=0.5pt,
              font=\scriptsize\sffamily},
  hdr/.style={font=\scriptsize\sffamily\bfseries},
  lab/.style={font=\scriptsize\sffamily, text=ln-muted},
  io/.style={->, thick, ln-accent},
  pcap/.style={font=\footnotesize\sffamily, anchor=north}]
  % ================= (a) passthrough
  \draw[dev] (0,0) rectangle (36,6) node[midway] {\iflangja{デバイス}{device}};
  \draw[band] (0,8.5) rectangle (22,29);
  \node[hdr] at (11,26.5) {\iflangja{ハイパーバイザ}{hypervisor}};
  \node[box, minimum width=18mm, minimum height=8mm] (vd) at (11,16.5)
    {\iflangja{VMM レベル\\ドライバ}{VMM-level\\driver}};
  \draw[->, densely dashed] (vd.south) -- (11,6);
  \draw[vm] (0,31) rectangle (36,55);
  \node[hdr] at (18,52.5) {\iflangja{ゲスト VM}{guest VM}};
  \draw[gos] (2,33) rectangle (34,49.5);
  \node[lab, text=black] at (18,47.2) {\iflangja{ゲスト OS}{guest OS}};
  \node[box, minimum width=24mm, minimum height=7mm] at (18,39.5)
    {\iflangja{デバイスドライバ}{device driver}};
  \draw[io] (29,36) -- (29,6);
  \node[lab, rotate=90, text=ln-accent] at (32.2,19) {\iflangja{直接アクセス}{direct access}};
  \node[pcap] at (18,-2) {\iflangja{(a) パススルー}{(a) passthrough}};
  % ================= (b) hypervisor-direct
  \begin{scope}[xshift=41mm]
  \draw[dev] (0,0) rectangle (36,6) node[midway] {\iflangja{デバイス}{device}};
  \draw[band] (0,8.5) rectangle (36,29);
  \node[hdr, anchor=west] at (1.2,26.5) {\iflangja{ハイパーバイザ}{hypervisor}};
  \node[box, minimum width=30mm, minimum height=4.3mm] (em) at (18,21.7)
    {\iflangja{デバイスエミュレーション}{device emulation}};
  \node[box, minimum width=30mm, minimum height=4.3mm] (st) at (18,16.6)
    {\iflangja{I/O スタック}{I/O stack}};
  \node[box, minimum width=30mm, minimum height=4.3mm] (dd) at (18,11.5)
    {\iflangja{デバイスドライバ}{device driver}};
  \draw[vm] (0,31) rectangle (36,55);
  \node[hdr] at (18,52.5) {\iflangja{ゲスト VM}{guest VM}};
  \draw[gos] (2,33) rectangle (34,49.5);
  \node[lab, text=black] at (18,47.2) {\iflangja{ゲスト OS}{guest OS}};
  \node[box, minimum width=26mm, minimum height=7mm] (gd) at (18,39.5)
    {\iflangja{エミュレートされた\\デバイスのドライバ}{driver for the\\emulated device}};
  \draw[io] (gd.south) -- (em.north);
  \node[lab, text=ln-accent, anchor=west] at (19,34.3) {\iflangja{トラップ}{trap}};
  \draw[io, shorten >=0pt] (dd.south) -- (18,6);
  \end{scope}
  \node[pcap] at (59,-2) {\iflangja{(b) ハイパーバイザ直接}{(b) hypervisor-direct}};
  % ================= (c) split-device driver
  \begin{scope}[xshift=82mm]
  \draw[dev] (0,0) rectangle (36,6) node[midway] {\iflangja{デバイス}{device}};
  \draw[band] (0,8.5) rectangle (36,29);
  \node[hdr, anchor=west] at (1.2,26.5) {\iflangja{ハイパーバイザ}{hypervisor}};
  \draw[vm] (0,31) rectangle (17,55);
  \node[hdr] at (8.5,52.5) {\iflangja{ゲスト VM}{guest VM}};
  \node[box, fill=ln-accent!12, minimum width=14mm, minimum height=8mm] (fe)
    at (8.5,44) {\iflangja{フロント\\エンド}{front-end\\driver}};
  \draw[svm] (19,31) rectangle (36,55);
  \node[hdr] at (27.5,52.5) {\iflangja{サービス VM}{service VM}};
  \node[box, fill=ln-accent!12, minimum width=14mm, minimum height=8mm] (be)
    at (27.5,44) {\iflangja{バック\\エンド}{back-end\\driver}};
  \node[box, minimum width=14mm, minimum height=6mm] (nd) at (27.5,35.5)
    {\iflangja{ドライバ}{driver}};
  \draw[<->, thick, ln-accent] (fe.east) -- (be.west);
  \draw[->] (be.south) -- (nd.north);
  \draw[io] (nd.south) -- (27.5,6);
  \end{scope}
  \node[pcap] at (100,-2) {\iflangja{(c) 分割デバイスドライバ}{(c) split-device driver}};
\end{tikzpicture}
   (width ~118 mm)
\begin{tikzpicture}[x=1mm, y=1mm, font=\scriptsize\sffamily,
  >={Stealth[length=1.5mm]},
  dev/.style={draw, fill=white},
  band/.style={draw, fill=ln-accent!22},
  vm/.style={draw=ln-rule, rounded corners=2pt, fill=ln-box},
  svm/.style={draw=ln-accent, rounded corners=2pt, fill=ln-accent!8},
  gos/.style={draw, fill=ln-accent!12},
  box/.style={draw, fill=white, align=center, inner sep=0.5pt,
              font=\scriptsize\sffamily},
  hdr/.style={font=\scriptsize\sffamily\bfseries},
  lab/.style={font=\scriptsize\sffamily, text=ln-muted},
  io/.style={->, thick, ln-accent},
  pcap/.style={font=\footnotesize\sffamily, anchor=north}]
  % ================= (a) passthrough
  \draw[dev] (0,0) rectangle (36,6) node[midway] {\iflangja{デバイス}{device}};
  \draw[band] (0,8.5) rectangle (22,29);
  \node[hdr] at (11,26.5) {\iflangja{ハイパーバイザ}{hypervisor}};
  \node[box, minimum width=18mm, minimum height=8mm] (vd) at (11,16.5)
    {\iflangja{VMM レベル\\ドライバ}{VMM-level\\driver}};
  \draw[->, densely dashed] (vd.south) -- (11,6);
  \draw[vm] (0,31) rectangle (36,55);
  \node[hdr] at (18,52.5) {\iflangja{ゲスト VM}{guest VM}};
  \draw[gos] (2,33) rectangle (34,49.5);
  \node[lab, text=black] at (18,47.2) {\iflangja{ゲスト OS}{guest OS}};
  \node[box, minimum width=24mm, minimum height=7mm] at (18,39.5)
    {\iflangja{デバイスドライバ}{device driver}};
  \draw[io] (29,36) -- (29,6);
  \node[lab, rotate=90, text=ln-accent] at (32.2,19) {\iflangja{直接アクセス}{direct access}};
  \node[pcap] at (18,-2) {\iflangja{(a) パススルー}{(a) passthrough}};
  % ================= (b) hypervisor-direct
  \begin{scope}[xshift=41mm]
  \draw[dev] (0,0) rectangle (36,6) node[midway] {\iflangja{デバイス}{device}};
  \draw[band] (0,8.5) rectangle (36,29);
  \node[hdr, anchor=west] at (1.2,26.5) {\iflangja{ハイパーバイザ}{hypervisor}};
  \node[box, minimum width=30mm, minimum height=4.3mm] (em) at (18,21.7)
    {\iflangja{デバイスエミュレーション}{device emulation}};
  \node[box, minimum width=30mm, minimum height=4.3mm] (st) at (18,16.6)
    {\iflangja{I/O スタック}{I/O stack}};
  \node[box, minimum width=30mm, minimum height=4.3mm] (dd) at (18,11.5)
    {\iflangja{デバイスドライバ}{device driver}};
  \draw[vm] (0,31) rectangle (36,55);
  \node[hdr] at (18,52.5) {\iflangja{ゲスト VM}{guest VM}};
  \draw[gos] (2,33) rectangle (34,49.5);
  \node[lab, text=black] at (18,47.2) {\iflangja{ゲスト OS}{guest OS}};
  \node[box, minimum width=26mm, minimum height=7mm] (gd) at (18,39.5)
    {\iflangja{エミュレートされた\\デバイスのドライバ}{driver for the\\emulated device}};
  \draw[io] (gd.south) -- (em.north);
  \node[lab, text=ln-accent, anchor=west] at (19,34.3) {\iflangja{トラップ}{trap}};
  \draw[io, shorten >=0pt] (dd.south) -- (18,6);
  \end{scope}
  \node[pcap] at (59,-2) {\iflangja{(b) ハイパーバイザ直接}{(b) hypervisor-direct}};
  % ================= (c) split-device driver
  \begin{scope}[xshift=82mm]
  \draw[dev] (0,0) rectangle (36,6) node[midway] {\iflangja{デバイス}{device}};
  \draw[band] (0,8.5) rectangle (36,29);
  \node[hdr, anchor=west] at (1.2,26.5) {\iflangja{ハイパーバイザ}{hypervisor}};
  \draw[vm] (0,31) rectangle (17,55);
  \node[hdr] at (8.5,52.5) {\iflangja{ゲスト VM}{guest VM}};
  \node[box, fill=ln-accent!12, minimum width=14mm, minimum height=8mm] (fe)
    at (8.5,44) {\iflangja{フロント\\エンド}{front-end\\driver}};
  \draw[svm] (19,31) rectangle (36,55);
  \node[hdr] at (27.5,52.5) {\iflangja{サービス VM}{service VM}};
  \node[box, fill=ln-accent!12, minimum width=14mm, minimum height=8mm] (be)
    at (27.5,44) {\iflangja{バック\\エンド}{back-end\\driver}};
  \node[box, minimum width=14mm, minimum height=6mm] (nd) at (27.5,35.5)
    {\iflangja{ドライバ}{driver}};
  \draw[<->, thick, ln-accent] (fe.east) -- (be.west);
  \draw[->] (be.south) -- (nd.north);
  \draw[io] (nd.south) -- (27.5,6);
  \end{scope}
  \node[pcap] at (100,-2) {\iflangja{(c) 分割デバイスドライバ}{(c) split-device driver}};
\end{tikzpicture}
   (width ~118 mm)
\begin{tikzpicture}[x=1mm, y=1mm, font=\scriptsize\sffamily,
  >={Stealth[length=1.5mm]},
  dev/.style={draw, fill=white},
  band/.style={draw, fill=ln-accent!22},
  vm/.style={draw=ln-rule, rounded corners=2pt, fill=ln-box},
  svm/.style={draw=ln-accent, rounded corners=2pt, fill=ln-accent!8},
  gos/.style={draw, fill=ln-accent!12},
  box/.style={draw, fill=white, align=center, inner sep=0.5pt,
              font=\scriptsize\sffamily},
  hdr/.style={font=\scriptsize\sffamily\bfseries},
  lab/.style={font=\scriptsize\sffamily, text=ln-muted},
  io/.style={->, thick, ln-accent},
  pcap/.style={font=\footnotesize\sffamily, anchor=north}]
  % ================= (a) passthrough
  \draw[dev] (0,0) rectangle (36,6) node[midway] {\iflangja{デバイス}{device}};
  \draw[band] (0,8.5) rectangle (22,29);
  \node[hdr] at (11,26.5) {\iflangja{ハイパーバイザ}{hypervisor}};
  \node[box, minimum width=18mm, minimum height=8mm] (vd) at (11,16.5)
    {\iflangja{VMM レベル\\ドライバ}{VMM-level\\driver}};
  \draw[->, densely dashed] (vd.south) -- (11,6);
  \draw[vm] (0,31) rectangle (36,55);
  \node[hdr] at (18,52.5) {\iflangja{ゲスト VM}{guest VM}};
  \draw[gos] (2,33) rectangle (34,49.5);
  \node[lab, text=black] at (18,47.2) {\iflangja{ゲスト OS}{guest OS}};
  \node[box, minimum width=24mm, minimum height=7mm] at (18,39.5)
    {\iflangja{デバイスドライバ}{device driver}};
  \draw[io] (29,36) -- (29,6);
  \node[lab, rotate=90, text=ln-accent] at (32.2,19) {\iflangja{直接アクセス}{direct access}};
  \node[pcap] at (18,-2) {\iflangja{(a) パススルー}{(a) passthrough}};
  % ================= (b) hypervisor-direct
  \begin{scope}[xshift=41mm]
  \draw[dev] (0,0) rectangle (36,6) node[midway] {\iflangja{デバイス}{device}};
  \draw[band] (0,8.5) rectangle (36,29);
  \node[hdr, anchor=west] at (1.2,26.5) {\iflangja{ハイパーバイザ}{hypervisor}};
  \node[box, minimum width=30mm, minimum height=4.3mm] (em) at (18,21.7)
    {\iflangja{デバイスエミュレーション}{device emulation}};
  \node[box, minimum width=30mm, minimum height=4.3mm] (st) at (18,16.6)
    {\iflangja{I/O スタック}{I/O stack}};
  \node[box, minimum width=30mm, minimum height=4.3mm] (dd) at (18,11.5)
    {\iflangja{デバイスドライバ}{device driver}};
  \draw[vm] (0,31) rectangle (36,55);
  \node[hdr] at (18,52.5) {\iflangja{ゲスト VM}{guest VM}};
  \draw[gos] (2,33) rectangle (34,49.5);
  \node[lab, text=black] at (18,47.2) {\iflangja{ゲスト OS}{guest OS}};
  \node[box, minimum width=26mm, minimum height=7mm] (gd) at (18,39.5)
    {\iflangja{エミュレートされた\\デバイスのドライバ}{driver for the\\emulated device}};
  \draw[io] (gd.south) -- (em.north);
  \node[lab, text=ln-accent, anchor=west] at (19,34.3) {\iflangja{トラップ}{trap}};
  \draw[io, shorten >=0pt] (dd.south) -- (18,6);
  \end{scope}
  \node[pcap] at (59,-2) {\iflangja{(b) ハイパーバイザ直接}{(b) hypervisor-direct}};
  % ================= (c) split-device driver
  \begin{scope}[xshift=82mm]
  \draw[dev] (0,0) rectangle (36,6) node[midway] {\iflangja{デバイス}{device}};
  \draw[band] (0,8.5) rectangle (36,29);
  \node[hdr, anchor=west] at (1.2,26.5) {\iflangja{ハイパーバイザ}{hypervisor}};
  \draw[vm] (0,31) rectangle (17,55);
  \node[hdr] at (8.5,52.5) {\iflangja{ゲスト VM}{guest VM}};
  \node[box, fill=ln-accent!12, minimum width=14mm, minimum height=8mm] (fe)
    at (8.5,44) {\iflangja{フロント\\エンド}{front-end\\driver}};
  \draw[svm] (19,31) rectangle (36,55);
  \node[hdr] at (27.5,52.5) {\iflangja{サービス VM}{service VM}};
  \node[box, fill=ln-accent!12, minimum width=14mm, minimum height=8mm] (be)
    at (27.5,44) {\iflangja{バック\\エンド}{back-end\\driver}};
  \node[box, minimum width=14mm, minimum height=6mm] (nd) at (27.5,35.5)
    {\iflangja{ドライバ}{driver}};
  \draw[<->, thick, ln-accent] (fe.east) -- (be.west);
  \draw[->] (be.south) -- (nd.north);
  \draw[io] (nd.south) -- (27.5,6);
  \end{scope}
  \node[pcap] at (100,-2) {\iflangja{(c) 分割デバイスドライバ}{(c) split-device driver}};
\end{tikzpicture}

  \caption{Device virtualization models.  (a) The VM's own driver accesses
    the device directly.  (b) The hypervisor intercepts device accesses and
    emulates them with its own I/O stack and driver.  (c) A front-end
    driver in the guest passes requests to a back-end driver in the service
    VM, which drives the device.}
  \label{fig:l12-devices}
\end{figure}

\needspace{6\baselineskip}
\subsection{Passthrough model}

\begin{definition}[Passthrough model]
  In the \term{passthrough model}[パススルーモデル] a VMM-level driver
  configures the access permissions of a device so that a VM accesses the
  device directly, with its own driver, bypassing the hypervisor.
\end{definition}

The VM obtains exclusive access to the device, and its device operations
follow the same path as on a physical machine.  This \emph{VMM bypass},
the counterpart of the operating system bypass of
Lesson~11,\index{operating system bypass} gives the best performance.  The
model has three drawbacks.
\begin{itemize}
  \item Sharing the device is difficult: it must be reassigned from one VM
    to another.
  \item The machine must have exactly the type of device that the driver in
    the VM expects, since that driver operates the hardware directly.
  \item Migration is tricky, because part of the state of the VM resides
    in the device, and the destination machine must provide an identical
    device into which that state can be transferred.
\end{itemize}

\subsection{Hypervisor-direct model}

\begin{definition}[Hypervisor-direct model]
  In the \term{hypervisor-direct model}[ハイパーバイザ直接モデル] the
  hypervisor intercepts every device access of a VM and emulates the device
  operation: it translates the operation into a generic I/O operation,
  passes it through the I/O stack that resides in the hypervisor, and
  invokes the hypervisor's own driver for the physical device.
\end{definition}

The key benefit is that the VM is decoupled from the physical device.  The
VM sees an emulated device, which can be the same on every machine,
independently of the physical device; the device can be shared among VMs;
migration becomes simpler; and the hypervisor deals with the specifics of
each device.  VMware ESX uses this model.  The costs are the latency that
interception and emulation add to every device operation, and the
complexity of the device driver ecosystem inside the hypervisor, which must
contain, and remain robust with, drivers for all supported devices.

\subsection{Split-device driver model}

\begin{definition}[Split-device driver model]
  In the \term{split-device driver model}[分割デバイスドライバモデル]
  device access is controlled by two cooperating drivers: a
  \term{front-end driver}[フロントエンドドライバ] in the guest VM, which
  presents the device interface to the guest OS, and a
  \term{back-end driver}[バックエンドドライバ] in the service VM or the
  host OS, which accesses the physical device through the regular driver.
\end{definition}

The front-end driver does not manipulate device registers that would have
to be trapped and emulated.  It passes complete requests, such as a packet
to send, to the back-end driver, in Xen through memory shared between the
two VMs \cite{barham2003}.  This eliminates the emulation overhead of the
hypervisor-direct model, and because all requests for a device pass
through the back-end driver, the service VM can manage how the device is
shared among the VMs.  The price is a modified guest: the front-end driver
is written for the hypervisor, which limits the model to paravirtualized
guests.  \Cref{tab:l12-devices} summarizes the three models.

\begin{table}[htbp]
  \centering
  \caption{Device virtualization models compared.}
  \label{tab:l12-devices}
  \small
  \begin{tabular}{@{}>{\raggedright}p{18mm}>{\raggedright}p{19mm}%
      >{\raggedright}p{20mm}>{\raggedright}p{23mm}%
      >{\raggedright\arraybackslash}p{23mm}@{}}
    \toprule
    Model & Physical device driven by & Driver in the guest & Sharing and
      migration & Cost per device operation \\
    \midrule
    Passthrough & guest VM & native driver & difficult & none beyond native \\
    \addlinespace
    Hypervisor-direct & hypervisor & driver for the emulated device
      & supported & interception and emulation \\
    \addlinespace
    Split-device driver & service VM or host OS & modified front-end
      & supported, managed by back-end & request passed to back-end \\
    \bottomrule
  \end{tabular}
\end{table}

\begin{keypoint}
  Devices lack a standard interface, so their virtualization trades
  performance against flexibility: passthrough is fastest but ties a VM to
  a device, hypervisor-direct decouples the VM at the cost of emulation,
  and the split-device driver model avoids emulation by modifying the
  guest's drivers.
\end{keypoint}

\section{Hardware support for virtualization}
\label{sec:l12-hw}

Around 2005, AMD and Intel added support for virtualization to their
processors: AMD with the technology code-named Pacifica, now AMD-V, and
Intel with Vanderpool, now Intel VT.\source{P3L6, hardware virtualization,
x86 VT revolution; Uhlig et al.\ \cite{uhlig2005}; Rosenblum and Garfinkel
\cite{rosenblum2005}.}

\begin{definition}[Hardware-assisted virtualization]
  Virtualization in which the processor itself provides the mechanisms
  needed to virtualize it, so that unmodified guest operating systems run
  under trap-and-emulate without binary translation or paravirtualization,
  is called \term{hardware-assisted virtualization}[ハードウェア支援仮想化].
\end{definition}

The extensions address the problems of the preceding sections.
\begin{description}
  \item[Closing the holes in the x86 ISA] In non-root mode the guest OS
    runs in ring~0, where instructions such as \texttt{POPF} behave as it
    expects, and every operation that the hypervisor must control, such as
    sensitive instructions, access to control registers, I/O and
    interrupts, can be made to cause a VM exit.  The architecture thereby
    becomes virtualizable by trap-and-emulate.
  \item[Root and non-root modes] The modes of \cref{sec:l12-rings}, called
    host and guest mode by AMD, let the guest OS run in ring~0 while the
    hypervisor keeps control.
  \item[VM control structure] A data structure per VM, more precisely per
    virtual CPU, that the hardware reads and updates.  It holds the state
    of the virtual CPU, saved on a VM exit and restored on a VM entry, and
    it specifies which instructions and events cause VM exits, so the
    hypervisor decides for each VM what to trap.
  \item[Extended page tables] The MMU performs the two-step translation of
    \cref{fig:l12-memory} itself.  The guest page table maps virtual to
    physical addresses as before, and an
    \term{extended page table}[拡張ページテーブル]
    (EPT)\index{EPT|see{extended page table}}, maintained by the
    hypervisor, maps physical to machine addresses.  On a TLB miss the MMU
    walks both tables, and the TLB caches the resulting translation from
    virtual to machine address.  The guest OS changes and switches its page
    tables without VM exits, and shadow page tables become unnecessary.
    AMD's equivalent is called nested page tables.
  \item[Tagged TLB] A \term{tagged TLB}[タグ付きTLB] marks each entry with
    an identifier of the VM to which it belongs.  Entries of different VMs
    coexist in the TLB, so a switch between VMs does not require flushing
    it.
  \item[Multiqueue devices and interrupt routing] Devices such as network
    cards provide several logical interfaces, or queues, each of which can
    be assigned to a different VM, and interrupts are routed to the core on
    which the VM that owns the queue runs.
  \item[Security and management support] Further extensions protect the
    hypervisor and support the management of VMs.
\end{description}
These features are used through new instructions, for example to enter
non-root mode, to read and write the VM control structure, and to issue
hypercalls.\margin{Intel's instructions include \texttt{VMLAUNCH} and
\texttt{VMRESUME} for VM entries, \texttt{VMREAD} and \texttt{VMWRITE} for
the control structure, and \texttt{VMCALL} for hypercalls.}  Hosted
hypervisors such as KVM rely on them.

Intel extended its support in stages, from the processor outwards.
\emph{VT-x} covers the processor: root and non-root modes, VM exits and
entries, the VM control structure, and, in later processors, extended page
tables and tagged TLB entries.  \emph{VT-d} (directed I/O) adds to the
chipset an I/O memory management unit that translates the addresses used
in DMA (Lesson~11) and remaps interrupts, so that a device can be assigned
safely to a VM in the passthrough model: the device cannot access the
memory of other VMs.  \emph{VT-c} (connectivity) brings virtualization
support to network devices, with multiple queues per device and devices
that present themselves as several virtual devices, each assignable to a
different VM.

\begin{table}[htbp]
  \centering
  \caption{Approaches to processor virtualization on x86.}
  \label{tab:l12-cpu}
  \small
  \begin{tabular}{@{}>{\raggedright}p{24mm}>{\raggedright}p{16mm}%
      >{\raggedright}p{27mm}>{\raggedright}p{22mm}%
      >{\raggedright\arraybackslash}p{14mm}@{}}
    \toprule
    Approach & Guest OS & Problematic instructions & Main overhead
      & Example \\
    \midrule
    Trap-and-emulate alone, before 2005 & unmodified & some fail silently:
      incorrect & traps & --- \\
    \addlinespace
    Binary translation & unmodified & rewritten at run time
      & translation, emulation & VMware \\
    \addlinespace
    Paravirtualization & modified & replaced by hypercalls & hypercalls
      & Xen \\
    \addlinespace
    Hardware-assisted & unmodified & cause VM exits & VM exits & KVM \\
    \bottomrule
  \end{tabular}
\end{table}

\begin{keypoint}[Lesson summary]
  Virtualization runs several operating systems concurrently on one
  machine; a hypervisor gives each VM an efficient, isolated duplicate of
  the machine.  Bare-metal hypervisors run on the hardware, usually with a
  service VM that owns the drivers; hosted hypervisors extend a host OS.
  The CPU is virtualized by trap-and-emulate, which x86 processors before
  2005 did not fully support, hence binary translation for unmodified
  guests and paravirtualization with hypercalls for modified ones.  Memory
  virtualization adds machine addresses beneath physical addresses and is
  handled with shadow page tables or with validated, batched page table
  updates.  Devices are virtualized by the passthrough, hypervisor-direct
  or split-device driver model.  Hardware-assisted virtualization, with
  root and non-root modes, VM exits, extended page tables and tagged TLBs,
  removes most of the software overhead.
\end{keypoint}

\end{document}
